% !TeX program = xelatex
\documentclass[11pt]{article}
\newcommand{\notesdir}{../../lecture-notes}
\providecommand{\notesdir}{.}

% -----------------------------------------------------------------------------
% Shared preamble for the CO5059 topic notes, the final exam review and the
% cheat sheet's sibling documents.
% Documents outside this folder define \notesdir (path to this folder) first.
% -----------------------------------------------------------------------------

% Page layout and typography
\usepackage[a4paper,portrait,margin=1in]{geometry}
\usepackage{fontspec}
\usepackage{microtype}
\usepackage{parskip}
\setlength{\emergencystretch}{1.5em}% fewer overfull lines

% Mathematics and diagrams
\usepackage{amsmath,amssymb,amsthm,mathtools}
\usepackage{tikz-cd}
\usetikzlibrary{arrows.meta,positioning,shapes.geometric,shapes.arrows,
  shapes.multipart,calc,fit,backgrounds,matrix,patterns,decorations.pathreplacing}

% Reusable styles for flowcharts and concept diagrams.
\tikzset{
  concept/.style={
    draw=brandprimary,
    fill=brandsecondary!8,
    rounded corners,
    minimum width=3.2cm,
    minimum height=0.9cm,
    align=center,
    font=\small,
  },
  diagram arrow/.style={-{Latex[length=2.5mm]},thick,brandprimary},
  % tree nodes (e.g. Tree-LRU) and hardware blocks
  tnode/.style={draw,rectangle,font=\small\ttfamily,inner sep=3pt,
    minimum height=0.55cm,fill=white},
  box/.style={draw,rectangle,align=center,font=\small,inner sep=4pt,
    fill=white},
  % ALU drawn vertically (inputs on the top edge, output at the bottom):
  % \pic at (x,y) {valu}; inputs at (x-0.65,y), (x+0.65,y); output (x,y-1)
  valu/.pic={\draw[fill=white] (-1.1,0)--(-0.2,0)--(0,-0.3)--(0.2,0)--(1.1,0)
    --(0.5,-1)--(-0.5,-1)--cycle; \node[font=\scriptsize] at (0,-0.62) {ALU};},
}

% Tables, figures, code, and lists
\usepackage{graphicx}
\graphicspath{{\notesdir/}}
\usepackage{enumitem}
\usepackage{booktabs}
% array is pinned: a newer array.sty in the local texmf tree needs a newer
% LaTeX kernel than the installed one; the dated rollback works with both.
\usepackage{array}[=2024-06-01]
% P{w}: paragraph column set ragged right (no stretched word gaps in tables)
\newcolumntype{P}[1]{>{\raggedright\arraybackslash}p{#1}}
\usepackage{listings}

% Color and links
\usepackage[dvipsnames]{xcolor}
% Same palette as the other course notes (co5115):
% brandprimary = navy lines/arrows/links, brandsecondary!NN = light blue fills.
\definecolor{brandprimary}{HTML}{030391}
\definecolor{brandsecondary}{HTML}{1488D8}
\usepackage[
  colorlinks=true,
  linkcolor=brandprimary,
  citecolor=brandsecondary,
  urlcolor=brandsecondary
]{hyperref}

% Headers and footers
\usepackage{fancyhdr}
\setlength{\headheight}{26pt}
\pagestyle{fancy}
\fancyhf{}
\fancyhead[L]{\small\textbf{\coursecode}\\\coursename}
\fancyhead[R]{\small HCMUT}
\fancyfoot[L]{\small Semester \semester}
\fancyfoot[R]{\small Page \thepage}

% -----------------------------------------------------------------------------
% Course information -- change these values once per course
% -----------------------------------------------------------------------------
\newcommand{\universityname}{HCMC University of Technology}
\newcommand{\faculty}{Faculty of Computer Science and Engineering}
\newcommand{\coursecode}{CO5059}
\newcommand{\coursename}{Advanced System Architectures}
\newcommand{\semester}{261}
\newcommand{\instructor}{Assoc.\ Prof.\ Tran Ngoc Thinh}
\newcommand{\instructoremail}{}% leave empty to hide
\newcommand{\notetaker}{Phan Nhat Minh}
\newcommand{\notetakeremail}{minh.phannhat95@hcmut.edu.vn}
\newcommand{\emaillink}[1]{%
  \href{mailto:#1}{\textcolor{brandprimary}{\texttt{#1}}}}

% Definitions, results, examples, proofs and notes in coloured boxes.
\input{\notesdir/colored-boxes}

% -----------------------------------------------------------------------------
% Course-specific notation
% -----------------------------------------------------------------------------
\newcommand{\CPI}{\mathrm{CPI}}
\newcommand{\AMAT}{\mathrm{AMAT}}

\lstdefinestyle{plain}{
  basicstyle=\ttfamily\small,
  showstringspaces=false,
  columns=fullflexible,
  keepspaces=true,
  breaklines=true,
  frame=single,
  framerule=0.3pt,
  rulecolor=\color{black!30},
  xleftmargin=4pt,
  xrightmargin=4pt,
}
\lstdefinestyle{mips}{
  style=plain,
  keywordstyle=\bfseries,
  morekeywords={lw,sw,add,sub,and,or,slt,beq,bne,beqz,xor,nop,addi,addiu,addu,
    subu,li,la,jal,jr,jalr,sll,srl,sra,sllv,srlv,srav,andi,ori,xori,nor,slti,
    sltiu,sltu,lui,mult,multu,div,divu,mfhi,mflo,ll,sc,syscall,move,lh,lhu,lb,
    lbu,sh,sb,blez,bgtz,bltz,bgez,j,mul,ld,sd,dadd,daddu,daddi,daddui,dsub,
    dsubu,bnez,l.d,s.d,add.d,sub.d,mul.d,div.d},
}
\lstset{style=mips}

% Long identifiers (GRADE_REPORT.Section_identifier) may break after "_".
\makeatletter
\DeclareRobustCommand{\_}{\ifmmode\nfss@text{\textunderscore}\else
  \textunderscore\allowbreak\fi}
\makeatother

% Marks a correction of an error in the slides: \begin{remark}\slidefix{} ...
\newcommand{\slidefix}{\textbf{\textcolor{red!75!black}{Slide correction.}}}
% Marks material added from the textbook or papers, not in the slides.
\newcommand{\beyond}{\textbf{Beyond the slides.}}

% -----------------------------------------------------------------------------
% Textbook citations as footnotes (book page + PDF page of textbooks/*.pdf)
% -----------------------------------------------------------------------------
% \srcnote{text}: a footnote that also works inside the coloured boxes
% (mdframed/tcolorbox would otherwise keep it inside the box with a letter).
% In a box it prints the mark and queues the text; the queue is emptied after
% the box (hooks below), so the note lands at the foot of the page.
\makeatletter
\newtoks\srcnote@pending
\def\srcnote@mp{mpfootnote}
\newcommand{\srcnote}[1]{%
  \ifx\@mpfn\srcnote@mp
    \stepcounter{footnote}%
    \protected@xdef\@thefnmark{\thefootnote}%
    \@footnotemark
    \xdef\srcnote@n{\the\c@footnote}%
    \global\srcnote@pending\expandafter{\the\expandafter\srcnote@pending
      \expandafter\footnotetext\expandafter[\srcnote@n]{#1}}%
  \else\footnote{#1}\fi}
\newcommand{\srcnote@flush}{\the\srcnote@pending\global\srcnote@pending{}}
\@for\srcnote@env:={theorem,definition,proposition,prop,lemma,corollary,%
exercise,tbexercise,example,eg,remark,note,solution,explanation,noot}\do{%
  \AddToHook{env/\srcnote@env/after}{\srcnote@flush}}
\makeatother
% Hennessy & Patterson, Computer Architecture: A Quantitative Approach, 5th ed.
% (references/Computer Architecture - A Quantitative Approach.pdf).
% Chapters: PDF page = printed page + 29.  Appendix X: page X-n is PDF n + offset.
\makeatletter
\@namedef{HPoff@}{29}
\@namedef{HPoff@A}{523}\@namedef{HPoff@B}{579}\@namedef{HPoff@C}{647}
\@namedef{HPoff@D}{736}\@namedef{HPoff@E}{804}\@namedef{HPoff@F}{831}
\@namedef{HPoff@G}{949}\@namedef{HPoff@H}{984}\@namedef{HPoff@I}{1029}
\@namedef{HPoff@J}{1076}\@namedef{HPoff@K}{1150}
% \hppage{X}{n}: "p. X-n (PDF p. m)"; X empty for a chapter page.
\newcommand{\hppage}[2]{p.~\if\relax\detokenize{#1}\relax\else#1-\fi#2\
  (PDF~\the\numexpr#2+\csname HPoff@#1\endcsname\relax)}
% \hppages{X}{first}{last}
\newcommand{\hppages}[3]{pp.~\if\relax\detokenize{#1}\relax\else#1-\fi#2--%
  \if\relax\detokenize{#1}\relax\else#1-\fi#3\
  (PDF~\the\numexpr#2+\csname HPoff@#1\endcsname\relax--%
  \the\numexpr#3+\csname HPoff@#1\endcsname\relax)}
\makeatother
% Footnote citations:
%   \HP{location}{page}            chapter page          \HP[last]{loc}{first}
%   \HPA{X}{location}{page}        appendix X page X-n   \HPA[last]{X}{loc}{first}
\newcommand{\HP}[3][]{\srcnote{H\&P 5th ed., #2,
  \if\relax\detokenize{#1}\relax\hppage{}{#3}\else\hppages{}{#3}{#1}\fi.}}
\newcommand{\HPA}[4][]{\srcnote{H\&P 5th ed., #3,
  \if\relax\detokenize{#1}\relax\hppage{#2}{#4}\else\hppages{#2}{#4}{#1}\fi.}}

% Textbook exercise. No solution is written: the box cites the exercise and
% the book pages with the theory needed to solve it.
% \begin{tbexercise}[appendix letter]{H&P number}{page}{theory pointer}
%   e.g. \begin{tbexercise}{1.8}{60}{\S1.9, \hppages{}{46}{48}}
%        \begin{tbexercise}[B]{B.1}{60}{\S B.2, \hppages{B}{15}{22}}
\makeatletter
\newenvironment{tbexercise}[4][]
  {\gdef\tbx@note{\srcnote{H\&P 5th ed., Exercise #2, \hppage{#1}{#3}.}}%
   \gdef\tbx@theory{#4}%
   \begin{exercise}[H\&P #2]}
  {\par\hfill{\small\itshape Theory: \tbx@theory}\tbx@note
   \end{exercise}}
\makeatother
% Source note under a diagram, for checking it against the original:
% \diagramsource{slides file, slide N; cf. textbook figure}
\newcounter{diagram}
% (the counter is stepped inside the note's line, so no page break can fall
% between the diagram and its note)
\newcommand{\diagramsource}[1]{\par\nopagebreak
  {\centering\footnotesize\itshape\color{black!65}\leavevmode\refstepcounter{diagram}%
   Diagram~\thediagram\ -- source: #1\par}}

% Source line under a section heading.
\newcommand{\sources}[1]{\par\noindent{\small\itshape\sloppy Sources: #1\par}\medskip}

% Cover page + contents page.  Usage: \makecover{Big title}
\newcommand{\makecover}[1]{%
\pagenumbering{roman}
\begin{titlepage}
  \thispagestyle{empty}
  \renewcommand{\emaillink}[1]{\href{mailto:##1}{\textcolor{black}{\texttt{##1}}}}% no colour on the cover

  \begin{center}
    \includegraphics[width=0.62\textwidth]{assets/text_logo_3.png}

    \vspace{0.35cm}
    {\large\bfseries\faculty\par}
  \end{center}

  \vfill
  \begin{center}
    {\Huge\bfseries #1\par}
    \vspace{0.55cm}
    {\LARGE\bfseries\coursename\par}
    \vspace{0.2cm}
    {\large\coursecode\par}

    \vspace{0.65cm}
    \rule{0.38\textwidth}{0.6pt}
    \vspace{0.65cm}

    {\large\textbf{Instructor:} \instructor
      \ifx\instructoremail\empty\else\ {\small(\emaillink{\instructoremail})}\fi\par}
    \vspace{0.3cm}
    {\large\textbf{Notes by:} \notetaker\
      {\small(\emaillink{\notetakeremail})}\par}
    \vspace{0.3cm}
    {\large\textbf{Semester:} \semester\par}
  \end{center}

  \vfill
  \begin{center}
    \small Last updated: \today
  \end{center}
\end{titlepage}

% -----------------------------------------------------------------------------
% Table of contents -- kept on its own page
% -----------------------------------------------------------------------------
\clearpage
\setcounter{page}{2}
\thispagestyle{empty}
\hypersetup{linkcolor=black}
\setcounter{tocdepth}{2}
\tableofcontents
\hypersetup{linkcolor=brandprimary}
\clearpage
\pagenumbering{arabic}
}


\begin{document}
\sloppy
\makecover{Final Exam Review}

\section*{Exam scope}
\addcontentsline{toc}{section}{Exam scope}
One section per topic in the final exam scope. Each section recaps what to know,
lists the usual traps, works through the lecturer's review questions, then gives
new exam-style questions with full solutions. The scope below is from the
lecturer's review slides; two A4 sheets of notes are allowed.
\begin{center}
\small
\begin{tabular}{@{}p{3.2cm}p{9.2cm}l@{}}
\toprule
Topic & What is examined & Section \\
\midrule
Introduction & performance: time, throughput, instruction types, CPI; Amdahl's Law & \ref{sec:intro} \\
Instruction set principles & all & \ref{sec:isa} \\
Pipelining & all & \ref{sec:pipe} \\
Cache & organization (fully associative, set associative, direct mapped);
replacement (LRU, Clock, Tree-LRU, FIFO); write policy (write-through, write-back,
write allocate, no-write allocate); the three categories of misses; performance
(AMAT, hit rate, miss rate, multilevel caches) & \ref{sec:cache} \\
Virtual memory & address translation, page table, TLB; paging, segmentation;
replacement & \ref{sec:vm} \\
Multiprocessors & Flynn classification, UMA, NUMA; synchronization (spin lock,
barrier); superpipelining, superscalar, multithreading; data dependences, in-order
and out-of-order issue and completion, VLIW; cache coherence (directory-based,
snoopy, MSI, MESI) & \ref{sec:mp} \\
\bottomrule
\end{tabular}
\end{center}
\clearpage

% ============================================================================
\section{Introduction}
\label{sec:intro}

This section covers the quantitative tools from the introduction lecture: how to
measure time and throughput, the CPU equation, instruction types and CPI, and
Amdahl's Law.

% ----------------------------------------------------------------------------
\subsection{What to know}

\subsubsection{Time and throughput}
\begin{definition}[Latency and throughput]
\textbf{Latency} (execution time, response time) is the time to finish
\emph{one} task. \textbf{Throughput} (bandwidth) is the number of tasks
finished per unit time.
\end{definition}
\begin{itemize}[leftmargin=*]
  \item If a machine runs one task at a time, throughput $=1/\text{latency}$.
  With overlap (several tasks in flight: pipelining, multiple cores, more
  passengers per plane) throughput can rise while latency stays the same or
  even gets worse.
  \item \textbf{Elapsed (wall-clock) time} includes I/O and other programs.
  \textbf{CPU time} $=$ user CPU time $+$ system CPU time; it excludes waiting.
  \item Unix \texttt{time}: \texttt{90.7u 12.9s 2:39 65\%} means
  $(90.7+12.9)/159=65\%$ of the elapsed time was CPU time. With several
  threads on several cores the percentage can exceed 100\%.
\end{itemize}

\subsubsection{Performance and speedup}
\[
  \text{Performance}=\frac{1}{\text{ExTime}}\qquad
  \text{Speedup of A over B}=\frac{\text{Perf}_A}{\text{Perf}_B}
  =\frac{\text{ExTime}_B}{\text{ExTime}_A}=\frac{\text{old time}}{\text{new time}}
\]
``X is $n\%$ faster than Y'' means
$\dfrac{\text{ExTime}_Y}{\text{ExTime}_X}=1+\dfrac{n}{100}$, so
$n=100\,\dfrac{\text{ExTime}_Y-\text{ExTime}_X}{\text{ExTime}_X}$ (divide by
the time of the \emph{faster} machine).

\subsubsection{The CPU equation}
\begin{theorem}[CPU time]
\[
  T=\frac{\text{Instructions}}{\text{Program}}\times
    \frac{\text{Cycles}}{\text{Instruction}}\times
    \frac{\text{Seconds}}{\text{Cycle}}
   =I\times\CPI\times C=\frac{I\times\CPI}{f}.
\]
\end{theorem}
\begin{center}
\begin{tabular}{@{}lccc@{}}
\toprule
Affects $\dots$ & $I$ & $\CPI$ & $C$ \\
\midrule
Program, compiler, ISA & X & X & \\
Organization (CPU design) & & X & X \\
Technology (VLSI) & & & X \\
\bottomrule
\end{tabular}
\end{center}

\subsubsection{Instruction types and CPI}
\begin{prop}[CPI of an instruction mix]
With $I_i$ instructions of class $i$, each taking $\CPI_i$ cycles, and
frequency $F_i=I_i/I$:
\[
  \text{Cycles}=\sum_i \CPI_i\,I_i,\qquad
  \CPI=\frac{\text{Cycles}}{I}=\sum_i \CPI_i\,F_i,
\]
\[
  \text{fraction of time in class } i=\frac{\CPI_i\,F_i}{\CPI}.
\]
\end{prop}
\begin{remark}
CPI is an \emph{average}, so it depends on the mix. A change that removes cheap
instructions can \emph{raise} the CPI while lowering the time: always compare
total cycles or time, never CPI or instruction count on their own.
\end{remark}

\subsubsection{Amdahl's Law}
\begin{theorem}[Amdahl's Law]
If an enhancement speeds up a fraction $F$ of the \emph{original execution
time} by a factor $S$, and the rest is unchanged,
\[
  \text{ExTime}_{new}=\text{ExTime}_{old}\left[(1-F)+\frac FS\right],\qquad
  \text{Speedup}=\frac{1}{(1-F)+\dfrac FS}\;\le\;\frac{1}{1-F}.
\]
\end{theorem}
\begin{itemize}[leftmargin=*]
  \item \textbf{Several disjoint enhancements} (every $F_i$ a fraction of the
  original time):
  $\text{Speedup}=\dfrac{1}{(1-\sum F_i)+\sum F_i/S_i}$.
  \item \textbf{Parallel processors}: with parallel fraction $p$ on $N$
  processors, $\text{Speedup}=\dfrac{1}{(1-p)+p/N}\le\dfrac1{1-p}$.
  \item \textbf{Target speedup $T$}: $S=\dfrac{F}{1/T-(1-F)}$; impossible when
  $T\ge1/(1-F)$. Required fraction: $F=\dfrac{1-1/T}{1-1/S}$.
  \item \textbf{$F$ from an instruction mix}: $F=\CPI_i F_i/\CPI$, the share of
  \emph{time}, not of instructions.
  \item \textbf{$f'$ measured after the enhancement} (share of the \emph{new}
  time spent in enhanced mode): $\text{Speedup}=(1-f')+f'S$, and the original
  fraction is $F=f'S/\big((1-f')+f'S\big)$.
  \item \textbf{Enhancement that slows the clock} by a factor $k$:
  $\text{Speedup}=\dfrac{1}{\big[(1-F)+F/S\big]\,k}$.
\end{itemize}

% ----------------------------------------------------------------------------
\subsection{Common traps}
\begin{enumerate}[leftmargin=*]
  \item Speedup is old time over new time, never the reverse; a speedup below 1
  is a slowdown.
  \item ``$n\%$ faster'' divides by the faster time: 10~s vs.\ 8~s is 25\%
  faster, while 8~s vs.\ 10~s is 20\% \emph{less time}.
  \item Converting clock rate: 2.5~GHz is a 0.4~ns cycle. Keep $f$ in Hz and
  $T$ in seconds until the last step.
  \item A higher clock rate or a lower CPI on its own does not mean faster; check
  all three factors of $I\times\CPI\times C$.
  \item In Amdahl's Law, $F$ is a fraction of \emph{time}. Given a mix, convert
  with $\CPI_iF_i/\CPI$ first.
  \item A fraction measured on the \emph{enhanced} machine is not the $F$ of the
  formula; use $(1-f')+f'S$.
  \item Cross-check every Amdahl answer with the CPU equation when you can:
  same $I$ and clock $\Rightarrow$ speedup $=\CPI_{old}/\CPI_{new}$.
\end{enumerate}

% ----------------------------------------------------------------------------
\subsection{Lecturer's review questions}
From the lecturer's review slides (\texttt{Review\_final\_2026.pdf}), with worked answers.

\begin{exercise}[Lecturer 1]\label{l:perf1}
In a program's execution time, multiplication instructions take 60\% and division
instructions 30\%. Multiplication can be made up to $6\times$ faster and division up
to $12\times$ faster. To make the computer $2\times$ faster by improving only
\emph{one} component, which component is selected?
(a) multiplication; (b) division; (c) both a and b are correct; (d) both a and b are
wrong.
\end{exercise}
\begin{solution}
\textbf{(a)}. Multiplication: $\dfrac{1}{0.4+0.6/6}=\dfrac{1}{0.5}=2.0$, which
reaches the target. Division: $\dfrac{1}{0.7+0.3/12}=\dfrac{1}{0.725}=1.38$; even
$S\to\infty$ gives only $1/0.7=1.43<2$.
\end{solution}

\begin{exercise}[Lecturer 2]
If both multiplication and division are improved, the new computer is faster by:
(1) 4.000; (2) 4.175; (3) 4.444; (4) none of the above.
\end{exercise}
\begin{solution}
\textbf{(3)}: $\dfrac{1}{(1-0.6-0.3)+0.6/6+0.3/12}=\dfrac{1}{0.1+0.1+0.025}
=\dfrac{1}{0.225}=4.444$.
\end{solution}

\begin{exercise}[Lecturer 3]
A program of 4 million instructions runs on a 2.5~GHz CPU in 8~ms. The average CPI
is: (a) 5; (b) 10; (c) 20; (d) none of the above.
\end{exercise}
\begin{solution}
\textbf{(a)}: $\CPI=\dfrac{T\times f}{I}=\dfrac{8\times10^{-3}\times2.5\times10^9}
{4\times10^6}=\dfrac{2\times10^7}{4\times10^6}=5$.
\end{solution}

\begin{exercise}[Lecturer 4]
A new compiler produces 5 million instructions with CPI 2; the clock is unchanged.
The speedup is: (a) 1.5; (b) 4; (c) 2; (d) none of the above.
\end{exercise}
\begin{solution}
\textbf{(c)}: $T'=\dfrac{5\times10^6\times2}{2.5\times10^9}=4$~ms, speedup
$=8/4=2$. More instructions, but far fewer cycles ($10^7$ vs.\ $2\times10^7$).
\end{solution}

% ----------------------------------------------------------------------------
\subsection{Exam-style questions}
Solutions are in Section~\ref{sec:intro-sol}.

\begin{exercise}[Latency vs.\ throughput]\label{q:intro-lat}
Web server A handles one request at a time and takes 2~ms per request. Web
server B handles 4 requests at the same time, each taking 5~ms.
\begin{enumerate}[label=(\alph*),leftmargin=*]
  \item Which server has the lower latency, and by what factor?
  \item Find the throughput of each server in requests per second. Which is
  higher, and by what percentage?
  \item 3000 requests arrive together. How long does each server take to finish
  all of them?
  \item Which server would you pick for one interactive user, and which for a
  busy data centre? Why?
  \item A program run on B reports \texttt{12.0u 2.0s 0:07 200\%}. Explain how
  the CPU percentage can exceed 100\%.
\end{enumerate}
\end{exercise}

\begin{exercise}[Three design options]\label{q:intro-design}
A program executes $4\times10^9$ instructions with $\CPI=1.6$ on a 2~GHz
processor. The design team has two options:
\begin{itemize}[leftmargin=*]
  \item \textbf{Option A}: a deeper pipeline raises the clock to 2.5~GHz, but
  the CPI rises to 1.9.
  \item \textbf{Option B}: new instructions cut the instruction count by 15\%;
  the clock stays at 2~GHz but the CPI rises to 1.75.
\end{itemize}
\begin{enumerate}[label=(\alph*),leftmargin=*]
  \item Find the CPU time of the original design and of each option.
  \item Which option is better? Give its speedup over the original and say how
  many percent faster it is.
  \item What CPI would option A need to match option B?
  \item Which factors of the CPU equation does each option change, and which
  part of the design (compiler, ISA, organization, technology) is responsible?
\end{enumerate}
\end{exercise}

\begin{exercise}[Instruction classes]\label{q:intro-mix}
A 2~GHz processor runs a program of $5\times10^9$ instructions with this mix:
\begin{center}
\begin{tabular}{@{}lcc@{}}
\toprule
Class & Frequency & CPI \\
\midrule
ALU & 40\% & 1 \\
Load/Store & 30\% & 3 \\
Branch & 20\% & 2 \\
Floating point & 10\% & 6 \\
\bottomrule
\end{tabular}
\end{center}
\begin{enumerate}[label=(\alph*),leftmargin=*]
  \item Find the average CPI.
  \item Find the fraction of execution time spent in each class.
  \item Find the CPU time and the MIPS rating.
  \item A compiler optimization removes half of the branches and changes nothing
  else. Find the new instruction count, CPI and CPU time. Explain why the CPI
  goes \emph{up} although the program gets faster.
\end{enumerate}
\end{exercise}

\begin{exercise}[Choosing an enhancement]\label{q:intro-choose}
Use the original mix of Exercise~\ref{q:intro-mix}. Two improvements are
proposed:
\begin{itemize}[leftmargin=*]
  \item \textbf{E1}: a faster FP unit lowers the FP CPI from 6 to 2.
  \item \textbf{E2}: a better cache lowers the load/store CPI from 3 to 2.
\end{itemize}
\begin{enumerate}[label=(\alph*),leftmargin=*]
  \item For each, find $F$ and $S$ and the overall speedup with Amdahl's Law.
  \item Check both answers with the CPU equation.
  \item Only 10\% of the instructions are FP but 30\% are loads/stores. Explain
  why E1 still wins.
  \item Find the speedup with both enhancements, with Amdahl's Law and with the
  CPU equation.
\end{enumerate}
\end{exercise}

\begin{exercise}[Limits and targets]\label{q:intro-limit}
A program spends 60\% of its execution time in matrix multiplication.
\begin{enumerate}[label=(\alph*),leftmargin=*]
  \item An accelerator runs matrix multiplication $20\times$ faster. Find the
  overall speedup.
  \item What is the largest overall speedup any accelerator could give?
  \item How fast must the accelerator be for an overall speedup of 2.4?
  \item Can an overall speedup of 2.5 be reached? Explain.
  \item What fraction of the time would have to be accelerated by $20\times$ to
  reach an overall speedup of 4?
\end{enumerate}
\end{exercise}

\begin{exercise}[Fraction measured after the enhancement]\label{q:intro-after}
After a SIMD unit is added (it runs vectorizable code $10\times$ faster), a
profiler on the \emph{new} machine shows that 40\% of the execution time is
spent in SIMD code.
\begin{enumerate}[label=(\alph*),leftmargin=*]
  \item Find the overall speedup due to the SIMD unit.
  \item What fraction of the \emph{original} execution time was vectorizable?
  \item A student plugs $F=0.4$, $S=10$ into Amdahl's Law. What does the
  student get, and why is it wrong?
  \item The SIMD unit is improved to $20\times$. Find the speedup over the
  original machine (without SIMD).
\end{enumerate}
\end{exercise}

\begin{exercise}[Enhancement with a clock penalty]\label{q:intro-clock}
An enhancement speeds up 30\% of the execution time by $4\times$, but its extra
logic makes the clock cycle 5\% longer.
\begin{enumerate}[label=(\alph*),leftmargin=*]
  \item Find the overall speedup with and without the clock penalty.
  \item How much longer could the clock cycle become before the enhancement
  stops paying off?
\end{enumerate}
\end{exercise}

\begin{exercise}[Amdahl's Law for multiprocessors]\label{q:intro-par}
90\% of a program can run in parallel.
\begin{enumerate}[label=(\alph*),leftmargin=*]
  \item Find the speedup on 8 and on 64 processors.
  \item What is the speedup limit?
  \item How many processors are needed for a speedup of 8?
  \item Comment on the efficiency (speedup per processor) in (a).
\end{enumerate}
\end{exercise}

\begin{exercise}[True or false]\label{q:intro-tf}
Answer and give a one-line reason.
\begin{enumerate}[label=(\alph*),leftmargin=*]
  \item A processor with a higher clock rate is always faster.
  \item Of two code sequences for the same task, the one with the lower CPI is
  always faster.
  \item On a machine that runs one task at a time, halving the latency doubles
  the throughput.
  \item In Amdahl's Law, $F$ can be taken as the fraction of instructions that
  benefit.
  \item With enough money, any overall speedup can be reached by speeding up
  only part of a program.
  \item If X is 25\% faster than Y, then Y takes 25\% more time than X.
\end{enumerate}
\end{exercise}

% ----------------------------------------------------------------------------
\subsection{Solutions}
\label{sec:intro-sol}

\begin{noot}{Exercise~\ref{q:intro-lat}}
\begin{enumerate}[label=(\alph*),leftmargin=*]
  \item A: 2~ms, B: 5~ms. A has $5/2=2.5\times$ lower latency.
  \item A: $1/0.002=500$ requests/s. B: $4/0.005=800$ requests/s. B has
  $800/500=1.6\times$, i.e.\ 60\% higher throughput.
  \item A: $3000\times2\text{ ms}=6$~s. B: $3000/4=750$ rounds of 5~ms
  $=3.75$~s.
  \item One interactive user sees only latency, so A. A data centre with many
  users cares about requests served per second, so B. Latency and throughput can
  rank machines in opposite orders.
  \item CPU time $=12.0+2.0=14$~s over an elapsed time of 7~s gives
  $14/7=200\%$. CPU time is summed over all cores, so two threads running on two
  cores at the same time give 2~s of CPU time per second of wall-clock time.
\end{enumerate}
\end{noot}

\begin{noot}{Exercise~\ref{q:intro-design}}
\begin{enumerate}[label=(\alph*),leftmargin=*]
  \item Original: $T=\dfrac{4\times10^9\times1.6}{2\times10^9}=3.2$~s.\\
  A: $T_A=\dfrac{4\times10^9\times1.9}{2.5\times10^9}=3.04$~s.\\
  B: $I_B=0.85\times4\times10^9=3.4\times10^9$,
  $T_B=\dfrac{3.4\times10^9\times1.75}{2\times10^9}=2.975$~s.
  \item B is better. Speedup $=3.2/2.975=1.076$, so B is 7.6\% faster than the
  original (A gives $3.2/3.04=1.053$). B is also
  $100\,(3.04-2.975)/2.975=2.2\%$ faster than A.
  \item $\CPI_A=\dfrac{T_B\times f_A}{I}=\dfrac{2.975\times2.5\times10^9}{4\times10^9}=1.859$.
  \item A changes $C$ (shorter cycle) and CPI: organization. B changes $I$
  and CPI: ISA (and the compiler that uses the new instructions).
\end{enumerate}
\end{noot}

\begin{noot}{Exercise~\ref{q:intro-mix}}
\begin{enumerate}[label=(\alph*),leftmargin=*]
  \item $\CPI=0.4(1)+0.3(3)+0.2(2)+0.1(6)=0.4+0.9+0.4+0.6=2.3$.
  \item ALU $0.4/2.3=17.4\%$, load/store $0.9/2.3=39.1\%$, branch
  $0.4/2.3=17.4\%$, FP $0.6/2.3=26.1\%$.
  \item $T=\dfrac{5\times10^9\times2.3}{2\times10^9}=5.75$~s.
  $\text{MIPS}=\dfrac{2\times10^9}{2.3\times10^6}=869.6$.
  \item Branches drop from $1\times10^9$ to $0.5\times10^9$, so
  $I'=4.5\times10^9$. Cycles $=11.5\times10^9-0.5\times10^9\times2=10.5\times10^9$.
  $\CPI'=10.5/4.5=2.33$, $T'=10.5\times10^9/(2\times10^9)=5.25$~s, a speedup of
  $5.75/5.25=1.095$. The removed branches (CPI 2) were cheaper than the average
  (2.3), so the remaining mix is more expensive per instruction, but there are
  fewer cycles in total. (MIPS also falls, to 857, although the program is
  faster.)
\end{enumerate}
\end{noot}

\begin{noot}{Exercise~\ref{q:intro-choose}}
\begin{enumerate}[label=(\alph*),leftmargin=*]
  \item E1: $F=0.6/2.3=0.261$, $S=6/2=3$:
  $\text{Speedup}=\dfrac{1}{0.739+0.261/3}=\dfrac{1}{0.739+0.087}=1.21$.\\
  E2: $F=0.9/2.3=0.391$, $S=3/2=1.5$:
  $\text{Speedup}=\dfrac{1}{0.609+0.391/1.5}=\dfrac{1}{0.609+0.261}=1.15$.
  \item E1: $\CPI'=2.3-0.1(6-2)=1.9$, $2.3/1.9=1.21$.
  E2: $\CPI'=2.3-0.3(3-2)=2.0$, $2.3/2.0=1.15$. Both match.
  \item What matters is the time saved per instruction,
  $F_i\,(\CPI_{old}-\CPI_{new})$: E1 saves $0.1\times4=0.4$ cycles, E2 saves
  $0.3\times1=0.3$ cycles. FP instructions are rare but expensive, so they take
  26\% of the time.
  \item Amdahl: $\dfrac{1}{(1-0.261-0.391)+0.261/3+0.391/1.5}
  =\dfrac{1}{0.348+0.087+0.261}=\dfrac{1}{0.696}=1.44$.
  CPU equation: $\CPI'=2.3-0.4-0.3=1.6$, $2.3/1.6=1.44$.
\end{enumerate}
\end{noot}

\begin{noot}{Exercise~\ref{q:intro-limit}}
\begin{enumerate}[label=(\alph*),leftmargin=*]
  \item $\dfrac{1}{0.4+0.6/20}=\dfrac{1}{0.43}=2.33$.
  \item $S\to\infty$: $1/(1-0.6)=2.5$.
  \item $S=\dfrac{0.6}{1/2.4-0.4}=\dfrac{0.6}{0.01667}=36$.
  \item No: 2.5 is the limit. It would need $1/2.5=0.4=(1-F)+F/S$, i.e.\
  $F/S=0$, an infinitely fast accelerator.
  \item $F=\dfrac{1-1/4}{1-1/20}=\dfrac{0.75}{0.95}=0.789$: about 79\% of the time.
\end{enumerate}
\end{noot}

\begin{noot}{Exercise~\ref{q:intro-after}}
\begin{enumerate}[label=(\alph*),leftmargin=*]
  \item Take the new time as 1: 0.6 is non-SIMD and 0.4 is SIMD. Without the
  unit the SIMD part would take $0.4\times10=4$, so the old time is
  $0.6+4=4.6$. Speedup $=4.6$.
  \item $F=4/4.6=0.870$: 87\% of the original time.
  \item $1/(0.6+0.04)=1.56$. The 40\% is a share of the \emph{new}, shorter
  time; Amdahl's $F$ must be a share of the \emph{old} time. Check:
  $1/(0.130+0.870/10)=1/0.217=4.6$.
  \item $F=20/23$: $\dfrac{1}{3/23+(20/23)/20}=\dfrac{1}{4/23}=5.75$.
\end{enumerate}
\end{noot}

\begin{noot}{Exercise~\ref{q:intro-clock}}
\begin{enumerate}[label=(\alph*),leftmargin=*]
  \item Without the penalty: new time $=0.7+0.3/4=0.775$, speedup
  $1/0.775=1.29$. With it: $0.775\times1.05=0.814$, speedup $1.23$.
  \item It pays off while $0.775\,k<1$, i.e.\ $k<1/0.775=1.29$: the cycle can
  grow by up to about 29\%.
\end{enumerate}
\end{noot}

\begin{noot}{Exercise~\ref{q:intro-par}}
\begin{enumerate}[label=(\alph*),leftmargin=*]
  \item $N=8$: $\dfrac{1}{0.1+0.9/8}=\dfrac{1}{0.2125}=4.71$.
  $N=64$: $\dfrac{1}{0.1+0.9/64}=\dfrac{1}{0.1141}=8.77$.
  \item $1/0.1=10$, however many processors.
  \item $1/8=0.125=0.1+0.9/N\Rightarrow N=0.9/0.025=36$.
  \item Efficiency: $4.71/8=59\%$ on 8 processors but only $8.77/64=14\%$ on 64.
  The serial 10\% dominates as $N$ grows.
\end{enumerate}
\end{noot}

\begin{noot}{Exercise~\ref{q:intro-tf}}
\begin{enumerate}[label=(\alph*),leftmargin=*]
  \item False: time is $I\times\CPI/f$; a higher CPI (Exercise~\ref{q:intro-design},
  option A) or a larger $I$ can cancel the faster clock.
  \item False: the lower-CPI sequence may execute more instructions; compare
  total cycles.
  \item True: throughput $=1/\text{latency}$ when tasks do not overlap.
  \item False: $F$ is a fraction of execution time; convert with
  $\CPI_iF_i/\CPI$.
  \item False: the speedup is bounded by $1/(1-F)$.
  \item True: $\text{ExTime}_Y/\text{ExTime}_X=1.25$. (But X takes 20\% less
  time than Y.)
\end{enumerate}
\end{noot}


% ============================================================================
\section{Instruction Set Principles}
\label{sec:isa}

This section covers everything on instruction sets in the course: ISA styles,
instruction set design (operations, addressing modes, operand types, formats,
registers), CISC vs.\ RISC, and the MIPS case study (formats, encoding,
addressing, constants, procedures, atomic operations and the assembler).

% ----------------------------------------------------------------------------
\subsection{What to know}

\subsubsection{The instruction set}
\begin{itemize}[leftmargin=*]
  \item The \textbf{instruction set} is the set of all instructions a
  processor can perform: the \textbf{interface between software and
  hardware}.
  \item Execution cycle: \textbf{instruction fetch} $\to$ \textbf{decode} $\to$
  \textbf{operand fetch} $\to$ \textbf{execute} (update status) $\to$
  \textbf{write back} $\to$ next fetch.
\end{itemize}

\subsubsection{ISA styles: where the operands live}
\begin{center}
\small
\begin{tabular}{@{}p{2.4cm}p{3.3cm}p{3.6cm}p{5.1cm}@{}}
\toprule
Style & $C=A+B$ & Operands & Pros / cons \\
\midrule
Stack & \texttt{PUSH A; PUSH B; ADD; POP C} & implicit, top of stack &
(+) code density, simple hardware; (--) low parallelism, stack is a bottleneck \\
Accumulator & \texttt{LOAD A; ADD B; STORE C} & one implicit register (AC) + memory &
(+) easy compiler, few instructions; (--) very high memory traffic, variable CPI \\
Memory-memory & \texttt{ADD C, A, B} & all in memory &
(+) simple to design and understand; (--) bottleneck, heavy memory access \\
Register-memory & \texttt{LOAD R1,A; ADD R3,R1,B; STORE R3,C} & register or memory &
(+) fewer instructions than load-store; (--) operand access time varies \\
Load-store (reg-reg) & \texttt{LOAD R1,A; LOAD R2,B; ADD R3,R1,R2; STORE R3,C} &
registers only; only loads/stores touch memory &
(+) simple fixed encoding, easy to pipeline, registers reused; (--) more
instructions \\
\bottomrule
\end{tabular}
\end{center}
Historical machines: Motorola 6800 (2 GPRs, accumulator, 1974), DEC VAX (16,
register-memory/memory-memory, 1977), Intel 8086 (1, extended accumulator,
1978), Motorola 68000 (16, register-memory, 1980), Intel 80386 (32,
register-memory, 1985), PowerPC and DEC Alpha (32, load-store, 1992). Other
styles: high-level-language architectures (Burroughs B5000, failed for lack of
compilers) and RISC.

\subsubsection{Instruction set design}
Five aspects: \textbf{operation repertoire}, \textbf{addressing modes},
\textbf{data types}, \textbf{instruction format}, \textbf{registers}. A simple
instruction has an \emph{opcode}, \emph{source operand reference(s)} and a
\emph{result operand reference}.
\begin{itemize}[leftmargin=*]
  \item \textbf{Instruction types} (4): data processing (arithmetic, logic),
  data storage (load/store to memory), data movement (register/register, I/O,
  \texttt{MOV}), control (test and branch: \texttt{CMP}, \texttt{JMP}).
  \item \textbf{Operator types} (8): arithmetic \& logical, data transfer,
  control (branch, jump, call/return, trap), system (OS call, virtual memory),
  floating point, decimal, string, graphics. The first four are standard on
  every machine; the others may be primitives, a coprocessor, or synthesized by
  the compiler.
\end{itemize}

\subsubsection{Addressing modes}
\begin{center}
\small
\begin{tabular}{@{}lll@{}}
\toprule
Mode & Example & Operand / effective address \\
\midrule
Immediate & \texttt{ADD R0,\#10} & the constant in the instruction \\
Register & \texttt{ADD R0,R1} & \texttt{R1} \\
Direct (absolute) & \texttt{ADD R0,(100)} & \texttt{M[100]} \\
Register indirect & \texttt{ADD R0,(R1)} & \texttt{M[R1]} \\
Displacement & \texttt{LD R1,100(R2)} & \texttt{M[R2+100]} \\
Indexed & \texttt{ADD R3,(R1+R2)} & \texttt{M[R1+R2]} \\
Memory indirect & \texttt{load r1 = M[M[r2]]} & \texttt{M[M[R2]]} \\
PC-relative & \texttt{branch r1<r3, 1000} & target $=$ PC $+$ offset \\
Scaled & \texttt{load r1 = M[100+r3+r4*d]} & \texttt{M[100+R3+R4$\times$d]} \\
\bottomrule
\end{tabular}
\end{center}
\begin{itemize}[leftmargin=*]
  \item Classified by the \emph{source} of data (immediate, register, memory)
  and by the \emph{address calculation} (direct, indirect, indexed).
  \item SPEC89 on the VAX: \textbf{displacement} (32--55\%) and
  \textbf{immediate} (17--43\%) dominate; register deferred 3--24\%, scaled
  0--16\%, memory indirect 1--6\%. So a new ISA should support these two well.
  \item Most displacements fit in 12--16 bits; small constants are about
  50\% of operands, so a 16-bit immediate field covers most cases.
\end{itemize}

\subsubsection{Operand types and sizes}
Character: ASCII 8-bit (almost always), Unicode 16-bit. Integer (2's
complement): short 16-bit, long 32-bit. Floating point: single 32-bit, double
64-bit. Business: BCD (exact decimal fractions). DSP: fixed point, block
floating point. Graphics: RGBA/XYZW with 8-, 16-bit or single-precision
components.

\subsubsection{Instruction format and registers}
\begin{itemize}[leftmargin=*]
  \item Encoding affects \textbf{program size} and \textbf{decode
  complexity}, and must fit enough registers and addressing modes.
  \item \textbf{Variable length} (VAX, x86): opcode $+$ address specifier/field
  pairs; dense code, hard to decode. \textbf{Fixed length} (ARM, MIPS,
  PowerPC): easy to decode and pipeline, larger code. \textbf{Hybrid}: two
  fixed lengths (MIPS16, Thumb) for density with simple decode.
  \item A field for $n$ choices needs $\lceil\log_2 n\rceil$ bits; $k$-bit
  two's-complement field: $-2^{k-1}\ldots2^{k-1}-1$.
  \item \textbf{Expanding opcode}: leave unused codes of a short opcode as
  escapes that extend into operand fields of shorter instructions.
  \item \textbf{Registers}: fastest storage, flexible for the compiler, but
  cost more per bit; more registers need more bits per instruction and make the
  register file slower.
\end{itemize}

\subsubsection{CISC vs.\ RISC}
\begin{center}
\small
\begin{tabular}{@{}p{3cm}p{5.6cm}p{5.6cm}@{}}
\toprule
& CISC & RISC \\
\midrule
Goal & do more per instruction, fewer instructions & simple instructions, $\CPI\le1$ \\
Motivation & expensive memory (M6800: 16K RAM \$500); backward compatibility & pipelining, compiler synthesizes complex ops \\
Encoding & variable length & fixed length (32 bits), few formats \\
Memory access & ALU operands can be in memory & only load/store \\
Addressing & many modes (14 in MC68000, 25 in MC68020) & immediate, register indirect, displacement, indexed \\
Other & 0- to 3-address instructions (VAX) & delayed loads/branches, prefetch, speculation; more instructions executed \\
Examples & VAX, 680x0, x86 & MIPS, SPARC, PowerPC, Alpha, PA-RISC \\
\bottomrule
\end{tabular}
\end{center}
Case studies: 680x0 has 18 addressing modes and instructions of 2--10 bytes;
x86 has 12 modes and instructions of 1--12 bytes; HP PA-RISC has 7 modes and 12
formats, all 32 bits; Alpha has 4 modes and 7 formats, all 32 bits.

\subsubsection{MIPS}
\textbf{Design principles}: simplicity favors regularity (fixed size, few
formats, opcode always in the first 6 bits); smaller is faster (few
instructions, registers, modes); make the common case fast (register operands,
immediates); good design demands good compromises (three formats).

SPEC2006 instruction mix: integer programs are 35\% data transfer, 34\%
conditional branch, 16\% arithmetic, 12\% logical, 2\% jump; floating-point
programs are 48\% arithmetic, 36\% data transfer, 8\% branch, 4\% logical.

\begin{center}
\small
\begin{tabular}{@{}lcccccc@{}}
\toprule
Format & 31--26 & 25--21 & 20--16 & 15--11 & 10--6 & 5--0 \\
\midrule
R & op (0) & rs & rt & rd & shamt & funct \\
I & op & rs & rt & \multicolumn{3}{c}{immediate (16)} \\
J & op & \multicolumn{5}{c}{address (26)} \\
\bottomrule
\end{tabular}
\end{center}

\begin{center}
\footnotesize
\begin{tabular}{@{}llll@{}}
\toprule
Mode in MIPS & Instructions & Address / operand & Range \\
\midrule
Register & R-format & \texttt{R[rs]}, \texttt{R[rt]} & 32 registers \\
Immediate & \texttt{addi}, \texttt{ori}, \texttt{slti} & 16-bit constant & $-32768\ldots32767$ \\
Base (displacement) & \texttt{lw}, \texttt{sw}, \texttt{lb}, \dots & \texttt{R[rs]} $+$ SignExt(imm) & $\pm32$ KiB \\
PC-relative & \texttt{beq}, \texttt{bne} & PC$+4+$SignExt(imm)$\times4$ & $\pm2^{15}$ instructions \\
Pseudo-direct & \texttt{j}, \texttt{jal} & $\{(\text{PC}{+}4)[31{:}28],\ \text{addr},\ 00\}$ & same 256 MiB region \\
Register indirect & \texttt{jr}, \texttt{jalr} & \texttt{R[rs]} & any address \\
\bottomrule
\end{tabular}
\end{center}

\begin{center}
\small
\begin{tabular}{@{}llll@{}}
\toprule
R-format (op $=0$) & funct & I/J-format & opcode \\
\midrule
\texttt{sll}, \texttt{srl}, \texttt{sra} & 00, 02, 03 & \texttt{j}, \texttt{jal} & 02, 03 \\
\texttt{jr}, \texttt{jalr} & 08, 09 & \texttt{beq}, \texttt{bne}, \texttt{blez}, \texttt{bgtz} & 04, 05, 06, 07 \\
\texttt{add}, \texttt{addu} & 20, 21 & \texttt{addi}, \texttt{addiu}, \texttt{slti}, \texttt{sltiu} & 08, 09, 0a, 0b \\
\texttt{sub}, \texttt{subu} & 22, 23 & \texttt{andi}, \texttt{ori}, \texttt{xori}, \texttt{lui} & 0c, 0d, 0e, 0f \\
\texttt{and}, \texttt{or}, \texttt{xor}, \texttt{nor} & 24, 25, 26, 27 & \texttt{lb}, \texttt{lh}, \texttt{lw}, \texttt{lbu}, \texttt{lhu} & 20, 21, 23, 24, 25 \\
\texttt{slt}, \texttt{sltu} & 2a, 2b & \texttt{sb}, \texttt{sh}, \texttt{sw}; \texttt{ll}, \texttt{sc} & 28, 29, 2b; 30, 38 \\
\bottomrule
\end{tabular}
\end{center}
(All in hex.)

\begin{center}
\small
\begin{tabular}{@{}llll@{}}
\toprule
Name & Number & Use & Preserved across a call? \\
\midrule
\texttt{\$zero} & 0 & constant 0 & n/a \\
\texttt{\$at} & 1 & assembler temporary & no \\
\texttt{\$v0--\$v1} & 2--3 & results & no \\
\texttt{\$a0--\$a3} & 4--7 & arguments & no \\
\texttt{\$t0--\$t7}, \texttt{\$t8--\$t9} & 8--15, 24--25 & temporaries & no \\
\texttt{\$s0--\$s7} & 16--23 & saved & yes \\
\texttt{\$k0--\$k1} & 26--27 & OS kernel & no \\
\texttt{\$gp}, \texttt{\$sp}, \texttt{\$fp}, \texttt{\$ra} & 28--31 & pointers; return address & yes \\
\bottomrule
\end{tabular}
\end{center}
\begin{remark}
\slidefix{} The slide table marks \texttt{\$a0--\$a3} as preserved; the MIPS
reference sheet (H\&P green card) says they are \emph{not}. A callee may
overwrite its arguments, so a caller that still needs them must save them.
\end{remark}

\begin{itemize}[leftmargin=*]
  \item \textbf{R-format}: \texttt{rd $\leftarrow$ rs funct rt}. Shifts use
  \texttt{rt} as the source, \texttt{shamt} as the amount, \texttt{rs}$=0$.
  \texttt{sra} shifts in the sign bit. \texttt{jr rs}: register indirect.
  \item \textbf{I-format}: \texttt{rt} is the destination (or the value stored
  by \texttt{sw}). \textbf{Sign-extended}: \texttt{addi}, \texttt{slti},
  load/store offsets, branch offsets. \textbf{Zero-extended}: \texttt{andi},
  \texttt{ori}, \texttt{xori}.
  \item \texttt{add}, \texttt{addi}, \texttt{sub} trap on signed overflow;
  \texttt{addu}, \texttt{addiu}, \texttt{subu} never trap. \texttt{slt}
  compares signed, \texttt{sltu} unsigned.
  \item \texttt{lb}/\texttt{lh} sign-extend, \texttt{lbu}/\texttt{lhu}
  zero-extend. Words must be 4-byte aligned, halfwords 2-byte aligned.
  \item \texttt{mult}/\texttt{div} write \texttt{HI}/\texttt{LO} (product:
  HI$=$upper 32 bits, LO$=$lower; division: LO$=$quotient, HI$=$remainder);
  read them with \texttt{mfhi}/\texttt{mflo}.
  \item \textbf{Big-endian} (as taught): the most significant byte is at the
  lowest address. Little-endian puts the least significant byte there.
  \item \textbf{Synthesized operations}: no \texttt{MOV}
  (\texttt{add \$t1,\$t0,\$zero}), no \texttt{NOT} (\texttt{nor \$d,\$s,\$zero}),
  no \texttt{SUBI} (\texttt{addi} with a negative constant), no \texttt{SLA}
  (same bits as \texttt{sll}).
  \item \textbf{32-bit constant}: \texttt{lui \$t0,upper; ori \$t0,\$t0,lower}
  (\texttt{ori} because it zero-extends).
  \item \textbf{Atomic read-modify-write}: \texttt{ll} then \texttt{sc} to the
  same address. \texttt{sc} stores and sets \texttt{rt}$=1$ only if the word was
  not written since the \texttt{ll}; otherwise it stores nothing and sets
  \texttt{rt}$=0$, so the code retries.
  \item \textbf{Procedures}: \texttt{jal label} puts the return address in
  \texttt{\$ra} and jumps; \texttt{jr \$ra} returns. Convention: save return
  address, save registers, run, restore registers, restore return address,
  return. The stack grows down: \texttt{addi \$sp,\$sp,-4k} then
  \texttt{sw}, and the reverse on exit. A non-leaf procedure must save
  \texttt{\$ra}.
  \item \textbf{Assembler}: directives \texttt{.text}, \texttt{.data},
  \texttt{.globl}, \texttt{.word}, \texttt{.space}. Pseudo-instructions
  (\texttt{li}, \texttt{la}, \texttt{move}, \texttt{lw} of a label, \texttt{blt},
  \dots) expand to real instructions, using \texttt{\$at} as scratch. SPIM's
  data segment starts at \texttt{0x10010000}.
\end{itemize}

% ----------------------------------------------------------------------------
\subsection{Common traps}
\begin{enumerate}[leftmargin=*]
  \item Branch offsets count \emph{instructions} from PC$+4$, not bytes from
  the branch: offset $=(\text{target}-(\text{PC}+4))/4$.
  \item The \texttt{j} field is the target \emph{divided by 4} (26 bits); the
  top 4 bits come from PC$+4$.
  \item Shifts: the source is \texttt{rt}, not \texttt{rs}; \texttt{rs}$=0$.
  \item I-format: the destination is \texttt{rt}, written \emph{second} in the
  binary but \emph{first} in assembly.
  \item \texttt{ori}/\texttt{andi} zero-extend; \texttt{addi} sign-extends.
  \texttt{lui}+\texttt{addi} breaks when the low half is $\ge$\texttt{0x8000}.
  \item \texttt{lb} vs.\ \texttt{lbu}: a byte $\ge$\texttt{0x80} becomes
  negative with \texttt{lb}.
  \item Code size and memory traffic are different questions: count
  instruction bytes for code size, data accesses for traffic.
  \item Fewer instructions (CISC) does not mean faster: use $I\times\CPI\times C$.
\end{enumerate}

% ----------------------------------------------------------------------------
\subsection{Lecturer's review questions}

\begin{exercise}[Lecturer 1]
In the MIPS instruction set, \texttt{NOT A} can be replaced by:
(a) \texttt{A NAND 0}; (b) \texttt{A XOR 0}; (c) \texttt{A NOR 0}; (d) none of the
above.
\end{exercise}
\begin{solution}
\textbf{(c)}: $\overline{A+0}=\overline{A}$, i.e.\ \texttt{nor \$d,\$s,\$zero}.
\texttt{A NAND 0} is always all ones and \texttt{A XOR 0} $=A$ (and MIPS has no
\texttt{nand}).
\end{solution}

\begin{exercise}[Lecturer 2]
Microprocessor performance can be measured by: (a) MIPS; (b) throughput; (c) MFLOPS;
(d) all of the above.
\end{exercise}
\begin{solution}
\textbf{(d)}. All three are performance \emph{metrics}. Remember the caveat from
Section~\ref{sec:intro}: MIPS and MFLOPS can mislead, and execution time of real
programs is the only reliable measure.
\end{solution}

\begin{exercise}[Lecturer 3]
The MIPS architecture is a: (a) stack; (b) accumulator; (c) register/register;
(d) register/memory architecture.
\end{exercise}
\begin{solution}
\textbf{(c)}: a load-store (register-register) ISA; only \texttt{lw}/\texttt{sw}
(and friends) access memory.
\end{solution}

\begin{exercise}[Lecturer 4]
Find \texttt{\$t1} after \texttt{lui \$t1,0x1234; addi \$t1,\$t1,0x5678}:
(a) \texttt{0x00000000}; (b) \texttt{0x1234}; (c) \texttt{0x5678};
(d) \texttt{0x12345678}.
\end{exercise}
\begin{solution}
\textbf{(d)}. \texttt{lui} gives \texttt{0x12340000}; \texttt{0x5678} has bit 15 $=0$,
so sign extension gives \texttt{0x00005678} and the sum is \texttt{0x12345678}.
(With a low half $\ge$\texttt{0x8000}, \texttt{addi} would subtract
\texttt{0x10000}; see Exercise~\ref{q:isa-sign}(e).)
\end{solution}

\begin{exercise}[Lecturer 5]
The x86 architecture is: (a) big-endian; (b) most-significant byte at the highest
address of a word; (c) most-significant byte at the lowest address of a word;
(d) none of the above.
\end{exercise}
\begin{solution}
\textbf{(b)}: x86 is \emph{little-endian}; the least-significant byte is at the
lowest address, so the most-significant byte is at the highest. (c) describes
big-endian, the MIPS convention in this course.
\end{solution}

% ----------------------------------------------------------------------------
\subsection{Exam-style questions}
Solutions are in Section~\ref{sec:isa-sol}.

\begin{exercise}[Comparing ISA styles]\label{q:isa-styles}
Compile $E=(A-B)\times(C+D)$, where $A$--$E$ are in memory, for a stack,
accumulator, memory-memory (3-operand), register-memory (3-operand, last operand
may be in memory, as in \texttt{ADD R3,R1,B}) and load-store machine. Do not
overwrite $A$--$D$; use memory temporaries \texttt{T}, \texttt{T1},
\texttt{T2} where needed.
Encoding: opcode 8 bits, memory address 16 bits, register field 4 bits (16
registers); stack \texttt{ADD}/\texttt{SUB}/\texttt{MUL} have only an opcode.
\begin{enumerate}[label=(\alph*),leftmargin=*]
  \item Write the code and count the instructions.
  \item Find the code size in bytes.
  \item Count the data memory accesses (reads $+$ writes of $A$--$E$ and
  temporaries). Assuming 4-byte data, find the total memory traffic (code $+$
  data) per execution.
  \item The next statement is $F=A\times B$. Which styles can reuse work
  already done?
\end{enumerate}
\end{exercise}

\begin{exercise}[Addressing modes]\label{q:isa-modes}
Registers: \texttt{R1}$=100$, \texttt{R2}$=200$, \texttt{R3}$=4$. Memory
(word at each address): \texttt{M[100]}$=300$, \texttt{M[104]}$=50$,
\texttt{M[116]}$=9$, \texttt{M[200]}$=100$, \texttt{M[204]}$=7$,
\texttt{M[300]}$=40$. Find the value loaded into \texttt{R4} and name the mode:
\begin{center}
\small
\begin{tabular}{@{}llll@{}}
(a) \texttt{LOAD R4,\#100} & (b) \texttt{LOAD R4,R1} & (c) \texttt{LOAD R4,(100)} & (d) \texttt{LOAD R4,(R2)} \\
(e) \texttt{LOAD R4,4(R2)} & (f) \texttt{LOAD R4,(R1+R3)} & (g) \texttt{LOAD R4,@(R1)} & (h) \texttt{LOAD R4,0(R1)[R3]} \\
\end{tabular}
\end{center}
(\texttt{@} is memory indirect; (h) is scaled: \texttt{M[0+R1+R3$\times$d]} with $d=4$.)
\begin{enumerate}[label=(\alph*),leftmargin=*,start=9]
  \item Which of these modes does MIPS have? Write MIPS code for (c), (f) and
  (g) with \texttt{\$s1}$=$\texttt{R1}, \texttt{\$s3}$=$\texttt{R3},
  \texttt{\$s4}$=$\texttt{R4}.
  \item Using the SPEC89 data, explain why MIPS keeps only a few modes.
\end{enumerate}
\end{exercise}

\begin{exercise}[Instruction format design]\label{q:isa-format}
\leavevmode
\begin{enumerate}[label=(\alph*),leftmargin=*]
  \item A machine has 16-bit instructions and 16 registers. It needs 15
  instructions with three register operands, 14 with two, 31 with one and 16
  with none. Design an expanding-opcode encoding, or show it is impossible.
  \item MIPS is extended to 128 registers, keeping 32-bit instructions and a
  6-bit opcode. How wide can the I-format immediate be now, and what is its
  range? How many bits are left in the R-format for \texttt{shamt} and
  \texttt{funct}? What are the consequences?
  \item Give two reasons for and two against adding more registers.
  \item Why does a fixed-length encoding help pipelining, and how do MIPS16 and
  Thumb recover code density?
\end{enumerate}
\end{exercise}

\begin{exercise}[CISC vs.\ RISC]\label{q:isa-cisc}
A program is compiled for a CISC and a RISC machine:
\begin{center}
\begin{tabular}{@{}lcc@{}}
\toprule
& CISC & RISC \\
\midrule
Static instruction count & 50{,}000 & 80{,}000 \\
Average instruction length & 3.5 bytes & 4 bytes \\
Dynamic instruction count & $1.0\times10^9$ & $1.6\times10^9$ \\
CPI & 4.0 & 1.2 \\
Clock & 2.0 GHz & 2.5 GHz \\
\bottomrule
\end{tabular}
\end{center}
\begin{enumerate}[label=(\alph*),leftmargin=*]
  \item Find the code size on each machine.
  \item Find the CPU time on each machine and the speedup of RISC.
  \item A compressed RISC mode (like Thumb) needs 15\% more instructions;
  70\% of them are 2 bytes and the rest 4 bytes. Find the new code size.
  \item Explain why CISC made sense in the 1970s and why it was kept.
\end{enumerate}
\end{exercise}

\begin{exercise}[Encoding MIPS]\label{q:isa-enc}
Give the format, fields, binary and hex of:
(a) \texttt{slt \$t0,\$s1,\$s2};
(b) \texttt{srl \$t1,\$s3,4};
(c) \texttt{andi \$s2,\$t3,0x00F0};
(d) \texttt{lbu \$t2,3(\$a1)};
(e) \texttt{sw \$ra,4(\$sp)};
(f) \texttt{jal Proc}, with \texttt{Proc} at \texttt{0x00400200}.
\end{exercise}

\begin{exercise}[Decoding MIPS]\label{q:isa-dec}
The words below are at addresses \texttt{0x00400000}, \texttt{0x00400004},
\dots{} in order. Decode each one; for branches and jumps give the target
address.
\begin{center}
\texttt{0x02119822}\quad \texttt{0x2108FFFF}\quad \texttt{0xAFBF0004}\quad
\texttt{0x1520FFFC}\quad \texttt{0x0C100080}
\end{center}
\end{exercise}

\begin{exercise}[Branch and jump ranges]\label{q:isa-range}
\leavevmode
\begin{enumerate}[label=(\alph*),leftmargin=*]
  \item \texttt{beq \$s0,\$s1,Far} is at \texttt{0x00400100} and \texttt{Far}
  is at \texttt{0x00480000}. Can it be encoded? If not, rewrite it and give the
  encoding of the jump.
  \item Can a \texttt{j} at \texttt{0x00400100} reach \texttt{0x10000000}? If
  not, how do you get there?
  \item State the reach of \texttt{beq}, \texttt{j} and \texttt{jr}.
  \item Why are branch and jump offsets stored in words rather than bytes?
\end{enumerate}
\end{exercise}

\begin{exercise}[Bytes, signs and overflow]\label{q:isa-sign}
MIPS is big-endian. The word \texttt{0x8A0B7CFF} is stored at address
\texttt{0x2000}.
\begin{enumerate}[label=(\alph*),leftmargin=*]
  \item Give the byte at each address \texttt{0x2000}--\texttt{0x2003}. Find
  the 32-bit result of \texttt{lb} and \texttt{lbu} from \texttt{0x2000} and
  from \texttt{0x2003}, \texttt{lh} from \texttt{0x2000} and \texttt{lhu} from
  \texttt{0x2002}.
  \item What does \texttt{lb} from \texttt{0x2000} return on a little-endian
  machine? What happens with \texttt{lw} from \texttt{0x2002}?
  \item \texttt{\$t1}$=$\texttt{0x7FFFFFFF}. What do \texttt{addi \$t2,\$t1,1}
  and \texttt{addiu \$t2,\$t1,1} do?
  \item \texttt{\$s0}$=$\texttt{0xFFFFFFFF}, \texttt{\$s1}$=1$. Find
  \texttt{slt \$t0,\$s0,\$s1} and \texttt{sltu \$t0,\$s0,\$s1}.
  \item Load \texttt{0x8000F000} into \texttt{\$t0}. What goes wrong with
  \texttt{lui \$t0,0x8000; addi \$t0,\$t0,0xF000}?
\end{enumerate}
\end{exercise}

\begin{exercise}[Directives and pseudo-instructions]\label{q:isa-asm}
In SPIM the data segment starts at \texttt{0x10010000}.
\begin{lstlisting}
        .data
x:      .word  5
buf:    .space 6
y:      .word  0x11223344
        .text
        .globl main
main:   la    $t0, y
        lw    $t1, x
        li    $t2, 0x12345
        move  $a0, $t1
\end{lstlisting}
\begin{enumerate}[label=(\alph*),leftmargin=*]
  \item Give the addresses of \texttt{x}, \texttt{buf} and \texttt{y}.
  \item Expand each instruction of \texttt{main} into real MIPS instructions.
  \item Why is \texttt{\$at} reserved for the assembler?
\end{enumerate}
\end{exercise}

\begin{exercise}[C to MIPS and instruction counts]\label{q:isa-count}
\leavevmode
\begin{lstlisting}[language=C,keywords={int,char,while,if,return}]
int count(char *s, char c) {      // s in $a0, c in $a1
    int n = 0;
    while (*s != 0) { if (*s == c) n++; s++; }
    return n;                     // in $v0
}
\end{lstlisting}
\begin{enumerate}[label=(\alph*),leftmargin=*]
  \item Write \texttt{count} in MIPS. Does it need the stack?
  \item How many instructions run for \texttt{s = "abca"}, \texttt{c = 'a'}?
  \item ALU instructions take 1 cycle and loads, branches and jumps take 2.
  Find the cycles and the CPI.
  \item Move the loop test to the bottom to remove the \texttt{j}. Repeat (b)
  and (c) and give the speedup.
\end{enumerate}
\end{exercise}

\begin{exercise}[Recursion and the stack]\label{q:isa-fact}
Use the recursive \texttt{fact} from the course:
\begin{lstlisting}
fact: addi $sp, $sp, -8
      sw   $ra, 4($sp)
      sw   $a0, 0($sp)
      slti $t0, $a0, 1
      beq  $t0, $zero, L1
      addi $v0, $zero, 1
      addi $sp, $sp, 8
      jr   $ra
L1:   addi $a0, $a0, -1
      jal  fact
      lw   $a0, 0($sp)
      lw   $ra, 4($sp)
      addi $sp, $sp, 8
      mul  $v0, $a0, $v0
      jr   $ra
\end{lstlisting}
\begin{enumerate}[label=(\alph*),leftmargin=*]
  \item \texttt{main} calls \texttt{fact(3)} with \texttt{\$sp}$=$\texttt{0x7FFFEFFC}.
  What is the lowest value \texttt{\$sp} reaches, and what is stored in each
  frame?
  \item How many instructions of \texttt{fact} run in total?
  \item Give \texttt{\$v0} at each return.
  \item Why must \texttt{\$a0} and \texttt{\$ra} be saved? What happens if
  \texttt{\$ra} is not?
\end{enumerate}
\end{exercise}

\begin{exercise}[Atomic increment with LL/SC]\label{q:isa-llsc}
\leavevmode
\begin{lstlisting}
again: ll   $t0, 0($s0)
       addi $t1, $t0, 1
       sc   $t1, 0($s0)
       beq  $t1, $zero, again
\end{lstlisting}
\begin{enumerate}[label=(\alph*),leftmargin=*]
  \item What does the code do, and what is in \texttt{\$t0} at the end?
  \item Two processors run it at the same time on a counter equal to 5. Trace
  one possible interleaving. What is the final value? What could happen with
  \texttt{lw}/\texttt{sw} instead?
  \item Why should the code between \texttt{ll} and \texttt{sc} be short?
\end{enumerate}
\end{exercise}

\begin{exercise}[Short answers]\label{q:isa-short}
Answer in one or two sentences.
\begin{enumerate}[label=(\alph*),leftmargin=*]
  \item Why does MIPS have no \texttt{SUBI}, no \texttt{NOT} and no
  \texttt{SLA}?
  \item List three uses of \texttt{\$zero}.
  \item Why is the opcode always in bits 31--26 and \texttt{rs}, \texttt{rt}
  always in the same place?
  \item Why is a load-store ISA easier to pipeline than a register-memory one?
  \item Which instruction classes should be fast for SPEC2006 integer and for
  floating-point programs?
  \item Why does MIPS need \texttt{ll}/\texttt{sc} instead of a single atomic
  instruction?
  \item Name an operand type a business machine needs that a DSP does not,
  and the reverse.
  \item Which part of the CPU equation does a RISC ISA increase, and which does
  it reduce?
\end{enumerate}
\end{exercise}

% ----------------------------------------------------------------------------
\subsection{Solutions}
\label{sec:isa-sol}

\begin{noot}{Exercise~\ref{q:isa-styles}}
(a), (b) Instruction sizes: \texttt{PUSH}/\texttt{POP}, accumulator
instructions $8+16=24$ bits; stack ALU $8$; memory-memory $8+3\cdot16=56$;
register load/store $8+4+16=28$; register-memory ALU $8+4+4+16=32$;
register-register ALU $8+3\cdot4=20$.
\begin{center}
\footnotesize
\begin{tabular}{@{}p{1.9cm}p{6.1cm}cp{2.5cm}c@{}}
\toprule
Style & Code & Instr. & Code size & Data acc. \\
\midrule
Stack & \texttt{PUSH A; PUSH B; SUB; PUSH C; PUSH D; ADD; MUL; POP E} & 8 &
$5\cdot24+3\cdot8=144$ b $=18$ B & 5 \\
Accumulator & \texttt{LOAD C; ADD D; STORE T; LOAD A; SUB B; MUL T; STORE E} & 7 &
$7\cdot24=168$ b $=21$ B & 7 \\
Memory-memory & \texttt{SUB T1,A,B; ADD T2,C,D; MUL E,T1,T2} & 3 &
$3\cdot56=168$ b $=21$ B & 9 \\
Register-memory & \texttt{LOAD R1,A; SUB R1,R1,B; LOAD R2,C; ADD R2,R2,D; MUL R3,R1,R2; STORE R3,E} & 6 &
$2(28{+}32){+}20{+}28$ $=168$ b $=21$ B & 5 \\
Load-store & \texttt{LOAD R1,A; LOAD R2,B; LOAD R3,C; LOAD R4,D; SUB R1,R1,R2; ADD R3,R3,R4; MUL R1,R1,R3; STORE R1,E} & 8 &
$5\cdot28+3\cdot20=200$ b $=25$ B & 5 \\
\bottomrule
\end{tabular}
\end{center}
(c) Total traffic $=$ code $+$ $4\times$data accesses: stack $18+20=38$ B,
accumulator $21+28=49$ B, memory-memory $21+36=57$ B, register-memory
$21+20=41$ B, load-store $25+20=45$ B. The stack machine has the densest code;
memory-memory and accumulator machines pay with data traffic (temporaries go to
memory); load-store has the largest code but the simplest instructions.

(d) Only the register styles: \texttt{A} and \texttt{B} are still in
\texttt{R1}/\texttt{R2} on the load-store machine (if the compiler keeps
them) and \texttt{A} in \texttt{R1} on the register-memory one, so no reloads
are needed. Stack, accumulator and memory-memory machines must read \texttt{A}
and \texttt{B} from memory again.
\end{noot}

\begin{noot}{Exercise~\ref{q:isa-modes}}
(a) immediate: 100. (b) register: 100. (c) direct: \texttt{M[100]}$=300$.
(d) register indirect: \texttt{M[200]}$=100$. (e) displacement:
\texttt{M[204]}$=7$. (f) indexed: \texttt{M[104]}$=50$. (g) memory indirect:
\texttt{M[M[100]]}$=$\texttt{M[300]}$=40$. (h) scaled:
\texttt{M[100+4$\times$4]}$=$\texttt{M[116]}$=9$.

(i) MIPS has register, immediate and displacement (plus PC-relative and
pseudo-direct for control). Register indirect is displacement with 0; direct is
displacement from \texttt{\$zero} (addresses below 32 KiB only):
\begin{lstlisting}
lw   $s4, 100($zero)      # (c) direct
add  $t0, $s1, $s3        # (f) indexed: form the address first
lw   $s4, 0($t0)
lw   $t0, 0($s1)          # (g) memory indirect: two loads
lw   $s4, 0($t0)
\end{lstlisting}
(j) Displacement and immediate cover most uses (32--55\% and 17--43\%); memory
indirect ($\le6\%$) and scaled ($\le16\%$) are rare and can be synthesized
with one extra instruction. Fewer modes keep decode simple and the clock short
(smaller is faster, make the common case fast).
\end{noot}

\begin{noot}{Exercise~\ref{q:isa-format}}
\begin{enumerate}[label=(\alph*),leftmargin=*]
  \item Possible, with every code used. Three operands: 4-bit opcode
  \texttt{0000}--\texttt{1110} (15), \texttt{1111} is the escape. Two operands:
  \texttt{1111 cccc rrrr rrrr}, \texttt{cccc} $=$ \texttt{0000}--\texttt{1101}
  (14), escapes \texttt{1110}, \texttt{1111}. One operand:
  \texttt{1111 111x cccc rrrr}, 32 codes, use 31 and keep
  \texttt{1111 1111 1111} as the escape. None: \texttt{1111 1111 1111 cccc},
  16 codes. Check: $15\cdot2^{12}+14\cdot2^8+31\cdot2^4+16=61440+3584+496+16=65536=2^{16}$.
  \item Register fields are 7 bits. I-format: $32-6-7-7=12$ bits, range
  $-2048\ldots2047$. R-format: $32-6-3\cdot7=5$ bits for \texttt{shamt} and
  \texttt{funct} together. Consequences: offsets and constants need
  \texttt{lui}/extra instructions more often, branches reach only $\pm2^{11}$
  instructions, and there is no room for both a 5-bit \texttt{shamt} and a
  \texttt{funct}, so fewer R-type operations (or variable shifts only).
  \item For: fewer loads/stores (spills), more values kept close for the
  compiler. Against: more bits per instruction, a slower (larger) register file,
  more state to save on calls and context switches.
  \item Every instruction has the same length and field positions, so the next
  PC is known (PC$+4$) and decode and register read can start at once, in fixed
  stages. MIPS16/Thumb add a second, 16-bit fixed format for common
  instructions: denser code, still simple decode.
\end{enumerate}
\end{noot}

\begin{noot}{Exercise~\ref{q:isa-cisc}}
\begin{enumerate}[label=(\alph*),leftmargin=*]
  \item CISC $50{,}000\times3.5=175{,}000$ B; RISC $80{,}000\times4=320{,}000$ B
  ($1.83\times$ larger).
  \item CISC $\dfrac{10^9\times4.0}{2\times10^9}=2.0$ s; RISC
  $\dfrac{1.6\times10^9\times1.2}{2.5\times10^9}=0.768$ s. Speedup
  $2.0/0.768=2.6$: 60\% more instructions, but a much lower CPI and faster clock.
  \item $80{,}000\times1.15=92{,}000$ instructions, average
  $0.7\cdot2+0.3\cdot4=2.6$ B, size $239{,}200$ B: 25\% smaller than plain RISC,
  still larger than CISC.
  \item Memory was expensive (16 KiB RAM for \$500, a 40 MB disk \$55{,}000),
  so compact code mattered, and complex instructions narrowed the gap to
  high-level languages. Later CISC designs had to run old binaries, so the
  instruction sets were kept for backward compatibility.
\end{enumerate}
\end{noot}

\begin{noot}{Exercise~\ref{q:isa-enc}}
\begin{center}
\footnotesize
\setlength{\tabcolsep}{3pt}
\begin{tabular}{@{}lp{4.1cm}l@{}}
\toprule
Instruction & Fields & Binary $=$ hex \\
\midrule
(a) R & op 0, rs 17, rt 18, rd 8, sh 0, fn \texttt{2a} &
\texttt{000000 10001 10010 01000 00000 101010} $=$ \texttt{0x0232402A} \\
(b) R & op 0, rs 0, rt 19, rd 9, sh 4, fn 02 &
\texttt{000000 00000 10011 01001 00100 000010} $=$ \texttt{0x00134902} \\
(c) I & op \texttt{0c}, rs 11, rt 18, imm \texttt{0x00F0} &
\texttt{001100 01011 10010 0000000011110000} $=$ \texttt{0x317200F0} \\
(d) I & op 24, rs 5, rt 10, imm 3 &
\texttt{100100 00101 01010 0000000000000011} $=$ \texttt{0x90AA0003} \\
(e) I & op \texttt{2b}, rs 29, rt 31, imm 4 &
\texttt{101011 11101 11111 0000000000000100} $=$ \texttt{0xAFBF0004} \\
(f) J & op 03, addr $=$ \texttt{0x00400200}$/4=$ \texttt{0x0100080} &
\texttt{000011 00000100000000000010000000} $=$ \texttt{0x0C100080} \\
\bottomrule
\end{tabular}
\end{center}
\end{noot}

\begin{noot}{Exercise~\ref{q:isa-dec}}
\begin{itemize}[leftmargin=*]
  \item \texttt{0x02119822} $=$ \texttt{000000 10000 10001 10011 00000 100010}:
  R, funct \texttt{22} $\Rightarrow$ \texttt{sub \$s3,\$s0,\$s1}.
  \item \texttt{0x2108FFFF} $=$ \texttt{001000 01000 01000 1111111111111111}:
  \texttt{addi \$t0,\$t0,-1}.
  \item \texttt{0xAFBF0004} $=$ \texttt{101011 11101 11111 0000000000000100}:
  \texttt{sw \$ra,4(\$sp)}.
  \item \texttt{0x1520FFFC} $=$ \texttt{000101 01001 00000 1111111111111100}:
  \texttt{bne \$t1,\$zero,-4}. At \texttt{0x0040000C}: target
  $=$ \texttt{0x00400010}$-16=$ \texttt{0x00400000}.
  \item \texttt{0x0C100080} $=$ \texttt{000011} $+$ \texttt{0x0100080}:
  \texttt{jal}, target $=\{0000,\ \texttt{0x0100080}\times4\}=$ \texttt{0x00400200}.
\end{itemize}
\end{noot}

\begin{noot}{Exercise~\ref{q:isa-range}}
\begin{enumerate}[label=(\alph*),leftmargin=*]
  \item Offset $=(\texttt{0x00480000}-\texttt{0x00400104})/4=131{,}007>32{,}767$:
  it does not fit. Invert the condition and jump:
\begin{lstlisting}
0x00400100        bne  $s0, $s1, Skip   # offset 1
0x00400104        j    Far              # 0x08120000
0x00400108  Skip: ...
\end{lstlisting}
  \texttt{j} field $=$ \texttt{0x00480000}$/4=$ \texttt{0x0120000}, encoding
  \texttt{0x08120000}; the top 4 bits of PC$+4$ and of \texttt{Far} are both 0.
  \item No: the top 4 bits of PC$+4$ are \texttt{0000} but the target's are
  \texttt{0001}. Use \texttt{lui \$at,0x1000; jr \$at}.
  \item \texttt{beq}: $-2^{15}\ldots2^{15}-1$ instructions from PC$+4$, about
  $\pm128$ KiB. \texttt{j}: anywhere in the current 256 MiB region.
  \texttt{jr}: any 32-bit address.
  \item Instructions are word-aligned, so the low 2 bits of every target are 0;
  storing words gives $4\times$ the reach for free.
\end{enumerate}
\end{noot}

\begin{noot}{Exercise~\ref{q:isa-sign}}
\begin{enumerate}[label=(\alph*),leftmargin=*]
  \item \texttt{0x2000}: \texttt{8A}, \texttt{0x2001}: \texttt{0B},
  \texttt{0x2002}: \texttt{7C}, \texttt{0x2003}: \texttt{FF}.
  \texttt{lb 0x2000} $=$ \texttt{0xFFFFFF8A} ($-118$);
  \texttt{lbu 0x2000} $=$ \texttt{0x0000008A} (138);
  \texttt{lb 0x2003} $=$ \texttt{0xFFFFFFFF} ($-1$);
  \texttt{lbu 0x2003} $=$ \texttt{0x000000FF} (255);
  \texttt{lh 0x2000} $=$ \texttt{0xFFFF8A0B};
  \texttt{lhu 0x2002} $=$ \texttt{0x00007CFF}.
  \item Little-endian puts \texttt{FF} at \texttt{0x2000}, so \texttt{lb}
  returns \texttt{0xFFFFFFFF}. \texttt{lw} from \texttt{0x2002} is not
  word-aligned: address error exception.
  \item \texttt{addi}: $2^{31}-1+1$ overflows, so an exception is raised and
  \texttt{\$t2} is not written. \texttt{addiu}: \texttt{\$t2}$=$\texttt{0x80000000}, no exception.
  \item \texttt{slt}: $-1<1$, so 1. \texttt{sltu}: $4{,}294{,}967{,}295<1$ is false, so 0.
  \item \texttt{lui \$t0,0x8000; ori \$t0,\$t0,0xF000}. With \texttt{addi},
  \texttt{0xF000} is sign-extended to \texttt{0xFFFFF000}; adding it to
  \texttt{0x80000000} ($-2^{31}$) gives signed overflow (an exception), and even
  \texttt{addiu} would give \texttt{0x7FFFF000}.
\end{enumerate}
\end{noot}

\begin{noot}{Exercise~\ref{q:isa-asm}}
\begin{enumerate}[label=(\alph*),leftmargin=*]
  \item \texttt{x} $=$ \texttt{0x10010000} (4 B); \texttt{buf} $=$
  \texttt{0x10010004}--\texttt{0x10010009}; \texttt{.word} aligns to 4, so
  \texttt{y} $=$ \texttt{0x1001000C} (2 bytes of padding).
  \item One valid expansion:
\begin{lstlisting}
lui  $at, 0x1001          # la $t0, y
ori  $t0, $at, 0x000C
lui  $at, 0x1001          # lw $t1, x
lw   $t1, 0($at)
lui  $at, 0x0001          # li $t2, 0x12345
ori  $t2, $at, 0x2345
addu $a0, $zero, $t1      # move $a0, $t1
\end{lstlisting}
  \item Pseudo-instructions need a scratch register to build addresses and
  constants; if programs used \texttt{\$at}, the expansion would destroy their
  values.
\end{enumerate}
\end{noot}

\begin{noot}{Exercise~\ref{q:isa-count}}
(a) A leaf procedure that only uses \texttt{\$t} registers needs no stack:
\begin{lstlisting}
count: add  $v0, $zero, $zero   # n = 0
loop:  lbu  $t0, 0($a0)         # *s
       beq  $t0, $zero, done
       bne  $t0, $a1, next
       addi $v0, $v0, 1         # n++
next:  addi $a0, $a0, 1         # s++ (chars are bytes)
       j    loop
done:  jr   $ra
\end{lstlisting}
(b) Each character: \texttt{lbu}, \texttt{beq}, \texttt{bne}, \texttt{addi},
\texttt{j} $=5$, plus 1 on a match. Total $1+4\cdot5+2+2+1=26$ (2 matches,
final \texttt{lbu}+\texttt{beq}, \texttt{jr}).

(c) ALU 7 (\texttt{add}, 2 \texttt{addi} $n$, 4 \texttt{addi} $s$), loads 5,
branches 9, jumps 5. Cycles $=7+2(5+9+5)=45$, $\CPI=45/26=1.73$.

(d)
\begin{lstlisting}
count: add  $v0, $zero, $zero
       lbu  $t0, 0($a0)
       beq  $t0, $zero, done
loop:  bne  $t0, $a1, next
       addi $v0, $v0, 1
next:  addi $a0, $a0, 1
       lbu  $t0, 0($a0)
       bne  $t0, $zero, loop
done:  jr   $ra
\end{lstlisting}
Instructions $=3+4\cdot4+2+1=22$: ALU 7, loads 5, branches 9, jumps 1.
Cycles $=7+2(5+9+1)=37$, $\CPI=37/22=1.68$. Speedup $=45/37=1.22$.
\end{noot}

\begin{noot}{Exercise~\ref{q:isa-fact}}
\begin{enumerate}[label=(\alph*),leftmargin=*]
  \item Four calls (\texttt{n}$=3,2,1,0$), 8 bytes each: lowest
  \texttt{\$sp} $=$ \texttt{0x7FFFEFFC}$-32=$ \texttt{0x7FFFEFDC}. Each frame
  holds \texttt{\$ra} at \texttt{4(\$sp)} and \texttt{\$a0} at
  \texttt{0(\$sp)}: frame 1 has \texttt{main}'s return address and 3; frames
  2--4 have the address after \texttt{jal fact} and 2, 1, 0.
  \item A call with \texttt{n}$\ge1$ runs 5 (entry and test) $+2$
  (\texttt{addi}, \texttt{jal}) $+5$ (after return) $=12$; the base case runs
  $5+3=8$. Total $3\cdot12+8=44$.
  \item \texttt{fact(0)}$=1$, \texttt{fact(1)}$=1\cdot1=1$,
  \texttt{fact(2)}$=2$, \texttt{fact(3)}$=6$.
  \item \texttt{\$a0} is changed for the recursive call but \texttt{n} is
  needed for the multiply afterwards; \texttt{jal} overwrites \texttt{\$ra}.
  Without saving \texttt{\$ra}, every \texttt{jr \$ra} returns to the
  instruction after \texttt{jal fact}, so the function never returns to
  \texttt{main} (it loops).
\end{enumerate}
\end{noot}

\begin{noot}{Exercise~\ref{q:isa-llsc}}
\begin{enumerate}[label=(\alph*),leftmargin=*]
  \item An atomic fetch-and-increment of \texttt{M[\$s0]}: \texttt{ll} reads
  the value, \texttt{sc} writes value$+1$ only if nobody wrote the word
  meanwhile, and otherwise the loop retries. \texttt{\$t0} holds the old value.
  \item P1 \texttt{ll} (5), P2 \texttt{ll} (5), P1 \texttt{sc} 6 succeeds, P2
  \texttt{sc} fails (the word changed) and returns 0, P2 retries: \texttt{ll}
  (6), \texttt{sc} 7. Final value 7. With \texttt{lw}/\texttt{sw} both could
  read 5 and write 6: one increment is lost.
  \item The longer the window, the more likely another write (or an interrupt)
  makes \texttt{sc} fail, so the loop may retry many times.
\end{enumerate}
\end{noot}

\begin{noot}{Exercise~\ref{q:isa-short}}
\begin{enumerate}[label=(\alph*),leftmargin=*]
  \item \texttt{SUBI} is \texttt{addi} with a negative constant; \texttt{NOT}
  is \texttt{nor \$d,\$s,\$zero}; a left arithmetic shift produces the same bits
  as \texttt{sll}. Leaving them out keeps the instruction set small.
  \item \texttt{move} (\texttt{add \$d,\$s,\$zero}), clear a register,
  compare with zero (\texttt{beq \$t0,\$zero,L}), negate
  (\texttt{sub \$d,\$zero,\$s}), direct addressing (\texttt{lw \$t0,100(\$zero)}).
  \item Regularity: the hardware can read \texttt{rs}/\texttt{rt} from the
  register file while it decodes the opcode, so decode is simple and fast.
  \item Only loads and stores touch memory, so every instruction uses each
  stage the same way (one memory access at most, in a fixed stage) and has a
  predictable length and timing.
  \item Integer: data transfers (35\%) and conditional branches (34\%).
  FP: arithmetic (48\%) and data transfers (36\%).
  \item A read-modify-write in one instruction needs both a read and a write
  in one memory operation, which breaks the load-store pipeline; the pair keeps
  each instruction simple and still guarantees atomicity.
  \item Business: BCD for exact decimal money values. DSP: fixed point or block
  floating point.
  \item It raises the instruction count $I$ and lowers the CPI (and allows a
  shorter clock cycle $C$).
\end{enumerate}
\end{noot}

% ============================================================================
\section{Pipelining}
\label{sec:pipe}

Everything on pipelining: the idea, the 5-stage MIPS pipeline and its timing,
structural, data and control hazards, forwarding, compiler scheduling and branch
schemes.

% ----------------------------------------------------------------------------
\subsection{What to know}

\subsubsection{The idea}
\begin{itemize}[leftmargin=*]
  \item \textbf{Pipelining} overlaps the execution of several instructions, each in
  a different stage. It improves \textbf{throughput}, not the latency of one
  instruction (latency even grows a little from pipeline-register overhead).
  \item The pipeline rate is limited by the \textbf{slowest stage}; unbalanced
  stages and the time to fill and drain reduce the speedup. Ideal speedup $=$
  number of stages.
  \item Laundry: wash 30, dry 40, fold 20 min, 4 loads: sequential
  $4\times90=360$ min, pipelined $30+4\times40+20=210$ min.
\end{itemize}

\subsubsection{The 5-stage MIPS pipeline}
\begin{center}
\small
\begin{tabular}{@{}lll@{}}
\toprule
Stage & Work & Resource \\
\midrule
IF & fetch instruction at PC, PC $\leftarrow$ PC$+4$ & instruction memory \\
ID & decode, read registers, sign-extend immediate & register file (read) \\
EX & ALU operation or effective address & ALU \\
MEM & load or store & data memory \\
WB & write the result register & register file (write) \\
\bottomrule
\end{tabular}
\end{center}
Pipeline registers IF/ID, ID/EX, EX/MEM, MEM/WB hold each instruction's values
between stages. The register file is \textbf{written in the first half} and
\textbf{read in the second half} of a cycle. MIPS is easy to pipeline: all
instructions are 32 bits, few formats with fixed register fields, memory operands
only in loads and stores, aligned memory accesses.

\subsubsection{Timing}
\begin{itemize}[leftmargin=*]
  \item \textbf{Clocked pipeline}, $k$ stages, $n$ instructions, $s$ stall cycles:
  \[
    \text{Cycles}=k+(n-1)+s,\qquad t_{clk}=\max_i t_i+t_{latch},\qquad
    \text{Time}=\text{Cycles}\times t_{clk}.
  \]
  \item Non-pipelined time per instruction $=\sum_i t_i$ (single-cycle: the clock
  is the longest instruction).
  \item Speedup, balanced stages: $\dfrac{nk}{k+n-1}\to k$. With stalls:
  \[
    \CPI_{pipe}=1+\text{stall cycles per instruction},
  \]
  \[
    \text{Speedup}=\frac{\CPI_{unpipe}}{\CPI_{pipe}}\times\frac{t_{clk,unpipe}}{t_{clk,pipe}}
    =\frac{\text{depth}}{1+\text{stalls/instr}}.
  \]
  \item \textbf{Stages that hand over as soon as the next one is free} (the
  lecturer's figure, no common clock): the first instruction takes $\sum t_i$, each
  later one finishes $\max t_i$ later:
  $T(n)=\sum_i t_i+(n-1)\max_i t_i$. An instruction waits in front of the slowest
  stage for (slowest $-$ slowest earlier stage) longer than the instruction before it.
\end{itemize}

\subsubsection{Hazards}
\begin{definition}[Hazards]
A \textbf{hazard} prevents the next instruction from executing in its cycle.
\textbf{Structural}: two instructions need the same hardware at once.
\textbf{Data}: an instruction needs a result that is not ready.
\textbf{Control}: the next PC depends on a branch not yet resolved.
\end{definition}
\begin{itemize}[leftmargin=*]
  \item \textbf{Structural}: e.g.\ one memory port for instructions and data, so a
  load stalls a later fetch. Fix: split I/D memories (caches) or duplicate the
  resource. It may be accepted when rare and costly to remove.
  \item \textbf{Data dependences}, $i$ before $j$:
  \textbf{RAW} ($j$ reads what $i$ writes, true dependence),
  \textbf{WAR} ($j$ writes what $i$ reads, anti-dependence),
  \textbf{WAW} ($j$ writes what $i$ writes, output dependence).
  \textbf{RAR is not a hazard} (nothing changes). In the in-order 5-stage pipeline
  only RAW causes stalls: reads (ID) are early and writes (WB) late and in order.
  WAR and WAW matter in out-of-order machines (Section~\ref{sec:mp}).
\end{itemize}

\begin{center}
\small
\begin{tabular}{@{}p{5cm}lp{6.2cm}@{}}
\toprule
Producer $\to$ consumer at distance $d$ & Stalls & Why \\
\midrule
No forwarding & $\max(0,3-d)$ & consumer's ID must be in the same cycle as producer's WB or later \\
Forwarding, ALU $\to$ ALU & 0 & EX/MEM $\to$ EX ($d=1$), MEM/WB $\to$ EX ($d=2$) \\
Forwarding, load $\to$ use, $d=1$ & 1 & data ready only at the end of MEM \\
Forwarding, load $\to$ use, $d\ge2$ & 0 & MEM/WB $\to$ EX \\
Forwarding, load $\to$ store data & 0 & MEM/WB $\to$ MEM \\
\bottomrule
\end{tabular}
\end{center}
\begin{itemize}[leftmargin=*]
  \item \textbf{Forwarding} (bypassing, hardware): feed the ALU from the pipeline
  registers instead of waiting for WB. The \textbf{load-use} hazard remains:
  hardware \textbf{interlocks} detect it and stall 1 cycle.
  \item \textbf{Compiler scheduling}: reorder independent instructions so no load is
  followed directly by its use. Reordering must keep every RAW, WAR and WAW order;
  if it cannot, rename registers.
  \item \textbf{Drawing a diagram}: a stalled instruction waits in ID
  (\textit{st} after ID); the one behind waits in IF (\textit{st} after IF); the
  next instruction is fetched when the IF stage is free. Total cycles $=$ cycle of
  the last WB.
\end{itemize}

\subsubsection{Control hazards}
\begin{itemize}[leftmargin=*]
  \item Branch penalty $=$ (stage that resolves the branch) $-1$: resolved in MEM
  $\Rightarrow$ 3 cycles; with the zero test and target adder moved into ID
  $\Rightarrow$ 1 cycle.
  \item $\CPI=1+f_{branch}\times(\text{fraction penalized})\times\text{penalty}$.
  \item Four alternatives: (1) \textbf{stall/flush} until resolved; (2)
  \textbf{predict not taken}: keep fetching PC$+4$, squash if taken (only taken
  branches pay); (3) \textbf{predict taken}: no gain in the 5-stage MIPS since the
  target is known no earlier than the outcome; (4) \textbf{delayed branch}: the
  instruction in the delay slot always executes.
  \item Filling the delay slot: (A) from before the branch (best, always useful,
  the branch must not depend on it); (B) from the target (useful when taken; must
  be safe when not); (C) from the fall-through (useful when not taken; must be safe
  when taken). Compilers fill about 60\% of slots, about 80\% of those usefully.
  \item \textbf{Compiler techniques}: instruction scheduling, delay-slot filling,
  static branch prediction, loop unrolling. \textbf{Hardware techniques}:
  forwarding, interlocks, early branch resolution, dynamic branch prediction.
\end{itemize}

% ----------------------------------------------------------------------------
\subsection{Common traps}
\begin{enumerate}[leftmargin=*]
  \item $n$ instructions take $k+n-1$ cycles, not $kn$ and not $n$.
  \item The clock is the \emph{slowest} stage (plus latch overhead), not the average.
  \item Without forwarding a dependence at distance 1 costs 2 stalls (not 3),
  because of the split-cycle register file.
  \item With forwarding only load-use at distance 1 stalls.
  \item Stalls push every later instruction right: the instruction behind a stall
  waits in IF.
  \item When reordering, check WAR and WAW too: moving a load above an instruction
  that still reads the old register value is wrong.
  \item Forwarding is hardware; scheduling and delay slots are compiler techniques.
\end{enumerate}

% ----------------------------------------------------------------------------
\subsection{Lecturer's review questions}

\begin{exercise}[Lecturer 1]
Choose the FALSE statement about techniques to reduce hazards \emph{using
compilers}: (a) forwarding; (b) scheduling; (c) branch prediction; (d) branch delay
slots.
\end{exercise}
\begin{solution}
\textbf{(a)}: forwarding is a hardware datapath feature. Scheduling and delay-slot
filling are done by the compiler, and branch prediction can be static (the compiler
or ISA chooses the predicted direction).
\end{solution}

\begin{exercise}[Lecturer 2]
Which data dependence is not a data hazard? (a) RAR; (b) RAW; (c) WAR; (d) WAW.
\end{exercise}
\begin{solution}
\textbf{(a)}: two reads of the same register can happen in any order.
\end{solution}

\begin{exercise}[Lecturer 3]
How long do 100 tasks take on an ideal 6-stage pipeline with 2~ns per stage?
(a) 12~ns; (b) 200~ns; (c) 210~ns; (d) 212~ns.
\end{exercise}
\begin{solution}
\textbf{(c)}: $(6+100-1)\times2=105\times2=210$~ns. (12~ns is the latency of one
task; 200~ns ignores the fill.)
\end{solution}

\begin{exercise}[Lecturer 4 and 5]
A 5-stage pipeline has stage times IF 5, ID 4, EX 5, MEM 10, WB 4~ns.
(4) The total execution time of 3 consecutive instructions is: (a) 38; (b) 43;
(c) 48; (d) 53~ns.
(5) The stall time of the 4th instruction is: (a) 5; (b) 10; (c) 15~ns; (d) none of
the above.
\end{exercise}
\begin{solution}
The answer choices fit the model in which an instruction moves on as soon as the
next stage is free (no common clock).
\textbf{(4) (c) 48~ns}: the first instruction takes $5+4+5+10+4=28$~ns; MEM, the
slowest stage, accepts one instruction every 10~ns, so each later instruction
finishes 10~ns later: $28+2\times10=48$~ns.
\begin{center}
\small
\begin{tabular}{@{}lccccc@{}}
\toprule
Instr. & IF & ID & EX & wait & MEM, WB \\
\midrule
1 & 0--5 & 5--9 & 9--14 & -- & 14--24, 24--28 \\
2 & 5--10 & 10--14 & 14--19 & 5 & 24--34, 34--38 \\
3 & 10--15 & 15--19 & 19--24 & 10 & 34--44, 44--48 \\
4 & 15--20 & 20--24 & 24--29 & 15 & 44--54, 54--58 \\
\bottomrule
\end{tabular}
\end{center}
\textbf{(5) (c) 15~ns}: instructions arrive in front of MEM every 5~ns (the slowest
earlier stage) but MEM takes 10~ns, so each one waits 5~ns longer than the one
before: $0, 5, 10, 15$~ns.

With a common clock instead, the cycle would be 10~ns and 3 instructions would take
$(5+2)\times10=70$~ns, which is not among the choices.
\end{solution}

\begin{exercise}[Lecturer 6: pipeline diagrams]\label{l:pipe-code}
The code runs on the 5-stage pipeline IF, ID, EX, MEM, WB:
\begin{lstlisting}
(1) lw   $t0, 0($t2)
(2) lw   $t1, 4($t4)
(3) add  $t3, $t0, $t1
(4) sw   $t3, 8($t2)
(5) lw   $t0, 12($t4)
(6) add  $t3, $t0, $t1
(7) sw   $t3, 16($t2)
\end{lstlisting}
\begin{enumerate}[label=(\alph*),leftmargin=*]
  \item Without forwarding, fill in the pipeline table, writing STALL where needed.
  \item Repeat with forwarding.
  \item Reorder the instructions so they run on the forwarding pipeline with as few
  stalls as possible.
\end{enumerate}
\end{exercise}
\begin{solution}
Dependences: RAW (1)$\to$(3) \texttt{\$t0}, (2)$\to$(3) \texttt{\$t1},
(3)$\to$(4) \texttt{\$t3}, (5)$\to$(6) \texttt{\$t0}, (2)$\to$(6) \texttt{\$t1},
(6)$\to$(7) \texttt{\$t3}; WAR (3)$\to$(5) on \texttt{\$t0}, (4)$\to$(6) on
\texttt{\$t3}; WAW (1)$\to$(5), (3)$\to$(6). \textit{st} $=$ STALL.

\textbf{(a) No forwarding}: a consumer's ID (register read) must be in the cycle of
the producer's WB or later.
\begin{center}\scriptsize\setlength{\tabcolsep}{1.1pt}
\begin{tabular}{@{}lccccccccccccccccccc@{}}
\toprule
Instruction & 1 & 2 & 3 & 4 & 5 & 6 & 7 & 8 & 9 & 10 & 11 & 12 & 13 & 14 & 15 & 16 & 17 & 18 & 19 \\
\midrule
(1) \texttt{lw \$t0,0(\$t2)} & IF & ID & EX & MEM & WB &  &  &  &  &  &  &  &  &  &  &  &  &  &  \\
(2) \texttt{lw \$t1,4(\$t4)} &  & IF & ID & EX & MEM & WB &  &  &  &  &  &  &  &  &  &  &  &  &  \\
(3) \texttt{add \$t3,\$t0,\$t1} &  &  & IF & ID & \textit{st} & \textit{st} & EX & MEM & WB &  &  &  &  &  &  &  &  &  &  \\
(4) \texttt{sw \$t3,8(\$t2)} &  &  &  & IF & \textit{st} & \textit{st} & ID & \textit{st} & \textit{st} & EX & MEM & WB &  &  &  &  &  &  &  \\
(5) \texttt{lw \$t0,12(\$t4)} &  &  &  &  &  &  & IF & \textit{st} & \textit{st} & ID & EX & MEM & WB &  &  &  &  &  &  \\
(6) \texttt{add \$t3,\$t0,\$t1} &  &  &  &  &  &  &  &  &  & IF & ID & \textit{st} & \textit{st} & EX & MEM & WB &  &  &  \\
(7) \texttt{sw \$t3,16(\$t2)} &  &  &  &  &  &  &  &  &  &  & IF & \textit{st} & \textit{st} & ID & \textit{st} & \textit{st} & EX & MEM & WB \\
\bottomrule
\end{tabular}
\end{center}
(3) waits for \texttt{\$t1} (WB in 6); (4) waits for \texttt{\$t3} (WB in 9); (6)
waits for \texttt{\$t0} (WB in 13); (7) waits for \texttt{\$t3} (WB in 16).
\textbf{19 cycles, 8 stall cycles.}

\textbf{(b) Forwarding}: only the two load-use pairs (2)$\to$(3) and (5)$\to$(6)
stall, 1 cycle each; (3)$\to$(4) and (6)$\to$(7) are forwarded from EX/MEM.
\begin{center}\scriptsize\setlength{\tabcolsep}{1.6pt}
\begin{tabular}{@{}lccccccccccccc@{}}
\toprule
Instruction & 1 & 2 & 3 & 4 & 5 & 6 & 7 & 8 & 9 & 10 & 11 & 12 & 13 \\
\midrule
(1) \texttt{lw \$t0,0(\$t2)} & IF & ID & EX & MEM & WB &  &  &  &  &  &  &  &  \\
(2) \texttt{lw \$t1,4(\$t4)} &  & IF & ID & EX & MEM & WB &  &  &  &  &  &  &  \\
(3) \texttt{add \$t3,\$t0,\$t1} &  &  & IF & ID & \textit{st} & EX & MEM & WB &  &  &  &  &  \\
(4) \texttt{sw \$t3,8(\$t2)} &  &  &  & IF & \textit{st} & ID & EX & MEM & WB &  &  &  &  \\
(5) \texttt{lw \$t0,12(\$t4)} &  &  &  &  &  & IF & ID & EX & MEM & WB &  &  &  \\
(6) \texttt{add \$t3,\$t0,\$t1} &  &  &  &  &  &  & IF & ID & \textit{st} & EX & MEM & WB &  \\
(7) \texttt{sw \$t3,16(\$t2)} &  &  &  &  &  &  &  & IF & \textit{st} & ID & EX & MEM & WB \\
\bottomrule
\end{tabular}
\end{center}
\textbf{13 cycles, 2 stall cycles.}

\textbf{(c) Reordering.} Moving (5) up between (2) and (3) looks attractive, but
(5) overwrites \texttt{\$t0} before (3) reads the value loaded by (1): a WAR
violation that changes the result. Keeping all dependences, (3) must follow both
(1) and (2), and nothing independent can go between them, so one load-use stall
remains. Swapping (4) and (5) removes the second one:
\begin{center}\scriptsize\setlength{\tabcolsep}{1.6pt}
\begin{tabular}{@{}lcccccccccccc@{}}
\toprule
Instruction & 1 & 2 & 3 & 4 & 5 & 6 & 7 & 8 & 9 & 10 & 11 & 12 \\
\midrule
(1) \texttt{lw \$t0,0(\$t2)} & IF & ID & EX & MEM & WB &  &  &  &  &  &  &  \\
(2) \texttt{lw \$t1,4(\$t4)} &  & IF & ID & EX & MEM & WB &  &  &  &  &  &  \\
(3) \texttt{add \$t3,\$t0,\$t1} &  &  & IF & ID & \textit{st} & EX & MEM & WB &  &  &  &  \\
(5) \texttt{lw \$t0,12(\$t4)} &  &  &  & IF & \textit{st} & ID & EX & MEM & WB &  &  &  \\
(4) \texttt{sw \$t3,8(\$t2)} &  &  &  &  &  & IF & ID & EX & MEM & WB &  &  \\
(6) \texttt{add \$t3,\$t0,\$t1} &  &  &  &  &  &  & IF & ID & EX & MEM & WB &  \\
(7) \texttt{sw \$t3,16(\$t2)} &  &  &  &  &  &  &  & IF & ID & EX & MEM & WB \\
\bottomrule
\end{tabular}
\end{center}
Order (1) (2) (3) (5) (4) (6) (7): \textbf{12 cycles, 1 stall}. If a free register
may be used, rename \texttt{\$t0} in (5)--(6) to \texttt{\$t5} and \texttt{\$t3} in
(6)--(7) to \texttt{\$t6}; then (5) can move up and there are no stalls:
\begin{center}\scriptsize\setlength{\tabcolsep}{1.6pt}
\begin{tabular}{@{}lccccccccccc@{}}
\toprule
Instruction & 1 & 2 & 3 & 4 & 5 & 6 & 7 & 8 & 9 & 10 & 11 \\
\midrule
(1) \texttt{lw \$t0,0(\$t2)} & IF & ID & EX & MEM & WB &  &  &  &  &  &  \\
(2) \texttt{lw \$t1,4(\$t4)} &  & IF & ID & EX & MEM & WB &  &  &  &  &  \\
(5) \texttt{lw \$t5,12(\$t4)} &  &  & IF & ID & EX & MEM & WB &  &  &  &  \\
(3) \texttt{add \$t3,\$t0,\$t1} &  &  &  & IF & ID & EX & MEM & WB &  &  &  \\
(4) \texttt{sw \$t3,8(\$t2)} &  &  &  &  & IF & ID & EX & MEM & WB &  &  \\
(6) \texttt{add \$t6,\$t5,\$t1} &  &  &  &  &  & IF & ID & EX & MEM & WB &  \\
(7) \texttt{sw \$t6,16(\$t2)} &  &  &  &  &  &  & IF & ID & EX & MEM & WB \\
\bottomrule
\end{tabular}
\end{center}
\textbf{11 cycles, 0 stalls} ($=5+7-1$, the ideal).
\end{solution}

% ----------------------------------------------------------------------------
\subsection{Exam-style questions}
Solutions are in Section~\ref{sec:pipe-sol}.

\begin{exercise}[Clocked pipeline]\label{q:pipe-clock}
Use the stage times of Lecturer question 4 (5, 4, 5, 10, 4~ns) and a 1~ns
pipeline-register overhead.
\begin{enumerate}[label=(\alph*),leftmargin=*]
  \item Find the clock cycle and the time for 1000 instructions on the clocked
  pipeline, and on a non-pipelined machine that takes 28~ns per instruction.
  \item MEM is split into two 5~ns stages. Repeat (a) and give the new speedup.
  \item Why is the speedup in (a) far below 5?
\end{enumerate}
\end{exercise}

\begin{exercise}[Finding hazards]\label{q:pipe-hz}
\leavevmode
\begin{lstlisting}
(1) add $t1, $t2, $t3
(2) sub $t4, $t1, $t5
(3) lw  $t2, 0($t4)
(4) or  $t1, $t6, $t7
(5) sw  $t2, 4($t1)
\end{lstlisting}
\begin{enumerate}[label=(\alph*),leftmargin=*]
  \item List every RAW, WAR and WAW dependence.
  \item Which of them can cause stalls in the 5-stage pipeline? Why not the others?
  \item Draw the diagram without forwarding and with forwarding. Give the cycles.
\end{enumerate}
\end{exercise}

\begin{exercise}[Branch schemes]\label{q:pipe-br}
20\% of instructions are branches and 60\% of them are taken. All other
instructions have $\CPI=1$. Find the CPI for:
(a) stall, branch resolved in MEM;
(b) predict not taken, resolved in MEM;
(c) stall, resolved in ID;
(d) predict not taken, resolved in ID;
(e) one delay slot (resolved in ID) filled usefully 50\% of the time.
(f) Give the speedup of (e) over (a).
\end{exercise}

\begin{exercise}[Delay slots]\label{q:pipe-slot}
\leavevmode
\begin{lstlisting}
      add  $s1, $s2, $s3
      beq  $s4, $zero, L
      (delay slot)
      sub  $t0, $t1, $t2
      ...
L:    or   $t3, $t4, $t5
\end{lstlisting}
\begin{enumerate}[label=(\alph*),leftmargin=*]
  \item Which instruction should fill the slot, and why is it the best choice?
  \item If the branch were \texttt{beq \$s1,\$zero,L}, why could \texttt{add} not
  be used? What are the other two options and their conditions?
\end{enumerate}
\end{exercise}

% ----------------------------------------------------------------------------
\subsection{Solutions}
\label{sec:pipe-sol}

\begin{noot}{Exercise~\ref{q:pipe-clock}}
\begin{enumerate}[label=(\alph*),leftmargin=*]
  \item $t_{clk}=10+1=11$~ns. Pipelined: $(5+999)\times11=11{,}044$~ns.
  Non-pipelined: $1000\times28=28{,}000$~ns. Speedup $=2.54$.
  \item 6 stages, $t_{clk}=5+1=6$~ns: $(6+999)\times6=6{,}030$~ns, speedup
  $28{,}000/6{,}030=4.64$.
  \item The stages are unbalanced: the 10~ns MEM stage sets the clock, so the
  4--5~ns stages idle half of each cycle, and the latch overhead adds to every
  stage.
\end{enumerate}
\end{noot}

\begin{noot}{Exercise~\ref{q:pipe-hz}}
\begin{enumerate}[label=(\alph*),leftmargin=*]
  \item RAW: (1)$\to$(2) \texttt{\$t1}, (2)$\to$(3) \texttt{\$t4}, (3)$\to$(5)
  \texttt{\$t2}, (4)$\to$(5) \texttt{\$t1}. WAR: (1)$\to$(3) \texttt{\$t2},
  (2)$\to$(4) \texttt{\$t1}. WAW: (1)$\to$(4) \texttt{\$t1}.
  \item Only RAW: registers are read in ID and written in WB, in program order, so
  a later instruction can never write before an earlier one reads (WAR) or writes
  (WAW).
  \item Without forwarding (15 cycles, 6 stalls):
\begin{center}\scriptsize\setlength{\tabcolsep}{1.6pt}
\begin{tabular}{@{}lccccccccccccccc@{}}
\toprule
Instruction & 1 & 2 & 3 & 4 & 5 & 6 & 7 & 8 & 9 & 10 & 11 & 12 & 13 & 14 & 15 \\
\midrule
(1) \texttt{add \$t1,\$t2,\$t3} & IF & ID & EX & MEM & WB &  &  &  &  &  &  &  &  &  &  \\
(2) \texttt{sub \$t4,\$t1,\$t5} &  & IF & ID & \textit{st} & \textit{st} & EX & MEM & WB &  &  &  &  &  &  &  \\
(3) \texttt{lw \$t2,0(\$t4)} &  &  & IF & \textit{st} & \textit{st} & ID & \textit{st} & \textit{st} & EX & MEM & WB &  &  &  &  \\
(4) \texttt{or \$t1,\$t6,\$t7} &  &  &  &  &  & IF & \textit{st} & \textit{st} & ID & EX & MEM & WB &  &  &  \\
(5) \texttt{sw \$t2,4(\$t1)} &  &  &  &  &  &  &  &  & IF & ID & \textit{st} & \textit{st} & EX & MEM & WB \\
\bottomrule
\end{tabular}
\end{center}
  With forwarding no stall is needed: (2), (3) and (5) get \texttt{\$t1},
  \texttt{\$t4} and \texttt{\$t1} from EX/MEM, and the load (3) is two
  instructions before the store that uses \texttt{\$t2} (9 cycles):
\begin{center}\scriptsize\setlength{\tabcolsep}{1.6pt}
\begin{tabular}{@{}lccccccccc@{}}
\toprule
Instruction & 1 & 2 & 3 & 4 & 5 & 6 & 7 & 8 & 9 \\
\midrule
(1) \texttt{add \$t1,\$t2,\$t3} & IF & ID & EX & MEM & WB &  &  &  &  \\
(2) \texttt{sub \$t4,\$t1,\$t5} &  & IF & ID & EX & MEM & WB &  &  &  \\
(3) \texttt{lw \$t2,0(\$t4)} &  &  & IF & ID & EX & MEM & WB &  &  \\
(4) \texttt{or \$t1,\$t6,\$t7} &  &  &  & IF & ID & EX & MEM & WB &  \\
(5) \texttt{sw \$t2,4(\$t1)} &  &  &  &  & IF & ID & EX & MEM & WB \\
\bottomrule
\end{tabular}
\end{center}
\end{enumerate}
\end{noot}

\begin{noot}{Exercise~\ref{q:pipe-br}}
(a) $1+0.2\times3=1.6$. (b) $1+0.2\times0.6\times3=1.36$. (c) $1+0.2\times1=1.2$.
(d) $1+0.2\times0.6\times1=1.12$. (e) the slot is wasted half the time:
$1+0.2\times0.5\times1=1.1$. (f) $1.6/1.1=1.45$.
\end{noot}

\begin{noot}{Exercise~\ref{q:pipe-slot}}
\begin{enumerate}[label=(\alph*),leftmargin=*]
  \item \texttt{add \$s1,\$s2,\$s3} from before the branch (strategy A): the branch
  does not use \texttt{\$s1}, and the instruction must execute on both paths, so
  the slot is always useful.
  \item The branch would need the \texttt{add}'s result, so it cannot be moved past
  it. From the target (B): copy \texttt{or} into the slot; useful when taken, and
  \texttt{\$t3} must be unused (dead) on the fall-through path. From the
  fall-through (C): move \texttt{sub} into the slot; useful when not taken, and
  \texttt{\$t0} must be dead at \texttt{L}.
\end{enumerate}
\end{noot}


% ============================================================================
\section{Cache}
\label{sec:cache}

Cache organization, block replacement (LRU, FIFO, Tree-LRU, Clock), write
policies, the three categories of misses, and cache performance (AMAT, hit and miss
rates, multilevel caches).

% ----------------------------------------------------------------------------
\subsection{What to know}

\subsubsection{Basics}
\begin{itemize}[leftmargin=*]
  \item Memory hierarchy: each level is larger, slower and cheaper per byte than
  the one above. It works because of \textbf{locality}: \emph{temporal} (an item
  used now is used again soon) and \emph{spatial} (nearby items are used soon, so
  whole blocks are moved).
  \item \textbf{Hit}: the data is in the cache. \textbf{Hit time}: time to access
  the cache, including deciding hit or miss. \textbf{Miss penalty}: time to bring
  the block from the lower level. Hit rate $+$ miss rate $=1$.
  \item The four questions: \emph{placement} (where can a block go),
  \emph{identification} (how is it found: tag $+$ valid bit), \emph{replacement}
  (which block is evicted), \emph{write strategy}.
\end{itemize}

\subsubsection{Organization and address breakdown}
$C$ = cache size, $B$ = block size, $a$ = associativity (ways), $A$ = address bits.
\[
  \#\text{blocks}=\frac CB,\qquad \#\text{sets } S=\frac{C}{B\times a},
\]
\[
  \text{offset}=\log_2B,\quad \text{index}=\log_2S,\quad \text{tag}=A-\text{index}-\text{offset}.
\]
\begin{center}
\small
\begin{tabular}{@{}lll@{}}
\toprule
Organization & Where a block can go & Index bits \\
\midrule
Direct mapped ($a=1$) & exactly one block: (block address) mod (\#blocks) & $\log_2(C/B)$ \\
$n$-way set associative & any of the $n$ ways of set (block address) mod $S$ & $\log_2(C/(Bn))$ \\
Fully associative ($S=1$) & anywhere; compare all tags & 0 \\
\bottomrule
\end{tabular}
\end{center}
\begin{itemize}[leftmargin=*]
  \item Block address $=\lfloor\text{addr}/B\rfloor$; set $=$ block address mod $S$;
  tag $=\lfloor\text{block address}/S\rfloor$.
  \item Doubling associativity (same $C$, $B$): index $-1$ bit, tag $+1$ bit.
  \item Storage per block $=8B$ data bits $+$ tag $+$ valid ($+$ dirty for write-back)
  ($+$ replacement bits).
\end{itemize}

\subsubsection{Block replacement}
Only needed when $a>1$ (direct mapped has no choice).
\begin{itemize}[leftmargin=*]
  \item \textbf{Random}: pseudo-random way; simple, no state.
  \item \textbf{LRU}: evict the block unused for the longest time. Exact LRU
  keeps the full use order of the set ($\lceil\log_2 n!\rceil$ bits, e.g.\ 5 bits
  for 4 ways), which is costly for large $n$.
  \item \textbf{FIFO}: evict the block loaded earliest; hits do not change the
  order.
  \item \textbf{Tree-LRU} (tree pseudo-LRU): $n-1$ bits in a binary tree over the
  ways (3 bits for 4 ways: root $b_0$ chooses a pair, $b_1$ the way in pair
  $\{0,1\}$, $b_2$ the way in pair $\{2,3\}$). On every access, set the bits on the
  path to point \emph{away} from the accessed way. Victim: follow the bits from the
  root. It approximates LRU and can evict a block that true LRU would keep.
  \item \textbf{Clock} (second chance): ways in a circle, one reference bit each and a
  pointer. On an access set the bit to 1. To replace: if the pointed block has bit 1,
  clear it and advance; evict the first block with bit 0. The new block gets bit 1
  and the pointer moves past it.
\end{itemize}

\subsubsection{Write policies}
\begin{center}
\small
\begin{tabular}{@{}p{2.6cm}p{11.6cm}@{}}
\toprule
Write-through & on a write hit, update the cache \emph{and} the next level; no dirty
bit; a write buffer hides the stalls; memory is always up to date. \\
Write-back & update only the cache; a \textbf{dirty bit} marks modified blocks, which
are written to memory only when evicted; less memory traffic. \\
\midrule
Write allocate & on a write miss, fetch the block into the cache, then write it. \\
No-write allocate & on a write miss, write only to the next level; the cache is not
changed. \\
\bottomrule
\end{tabular}
\end{center}
Usual pairs: write-through $+$ no-write allocate; write-back $+$ write allocate.

\subsubsection{The three categories of misses (3 Cs)}
\begin{itemize}[leftmargin=*]
  \item \textbf{Compulsory} (cold start): the first access to a block.
  \item \textbf{Capacity}: the cache cannot hold all the blocks the program needs;
  the miss also happens in a fully associative LRU cache of the same size.
  \item \textbf{Conflict} (collision): too many blocks map to the same set; the
  access would hit in a fully associative cache of the same size.
  \item (A fourth C, \textbf{coherence}, appears in multiprocessors.)
\end{itemize}
\begin{center}
\small
\begin{tabular}{@{}p{4.2cm}p{4.6cm}p{5.4cm}@{}}
\toprule
Basic optimization & Reduces & Cost \\
\midrule
Larger block size & compulsory misses (spatial locality) & larger miss penalty; more conflict/capacity misses if too large \\
Larger cache & capacity misses & longer hit time, cost, power \\
Higher associativity & conflict misses & longer hit time \\
Multilevel caches & miss penalty & more hardware \\
Read misses before writes & miss penalty & check the write buffer \\
Avoid translation in indexing & hit time & needs VIPT, $C/a\le$ page size \\
\bottomrule
\end{tabular}
\end{center}
Rule of thumb: a direct-mapped cache of size $N$ misses about as often as a 2-way
cache of size $N/2$.

\subsubsection{Cache performance}
\[
  \text{AMAT}=\text{Hit time}+\text{Miss rate}\times\text{Miss penalty}
\]
\[
  \CPI=\CPI_{exec}+\frac{\text{memory accesses}}{\text{instruction}}\times\text{Miss rate}\times\text{Miss penalty},
  \qquad T=I\times\CPI\times C.
\]
Accesses per instruction $=1$ (fetch) $+$ fraction of loads/stores. Miss penalty
in \emph{cycles} $=$ penalty in ns / clock period.

\textbf{Multilevel caches} (L2 reduces the L1 \emph{miss penalty}):
\[
  \text{AMAT}=\text{HT}_{L1}+\text{MR}_{L1}\times\big(\text{HT}_{L2}+\text{MR}_{L2}\times\text{MP}_{L2}\big)
\]
\begin{itemize}[leftmargin=*]
  \item \textbf{Local} miss rate $=$ misses in this cache / accesses \emph{to this
  cache} (for L2: accesses $=$ L1 misses). \textbf{Global} miss rate $=$ misses in
  this cache / all CPU accesses $=\text{MR}_{L1}\times\text{MR}_{L2,local}$.
  \item Stalls per instruction $=\dfrac{\text{L1 misses}}{\text{instr}}\times
  \text{HT}_{L2}+\dfrac{\text{L2 misses}}{\text{instr}}\times\text{MP}_{L2}$.
\end{itemize}

% ----------------------------------------------------------------------------
\subsection{Common traps}
\begin{enumerate}[leftmargin=*]
  \item Sets $=C/(B\times a)$: divide by the associativity.
  \item Work with block addresses (addr$/B$), not byte addresses, when tracing.
  \item L2 \emph{local} miss rate is relative to L1 misses; multiply by
  $\text{MR}_{L1}$ for the global rate.
  \item Count the instruction fetch in memory accesses per instruction.
  \item Tree-LRU is not LRU: trace the bits, do not guess.
  \item In write-back, a miss that evicts a dirty block costs a write-back as well.
  \item A conflict miss is one that a fully associative cache of the same size
  would \emph{not} have had.
\end{enumerate}

% ----------------------------------------------------------------------------
\subsection{Lecturer's review questions}

\begin{exercise}[Lecturer 1]
The algorithm that removes contents and places new contents into the cache is
called: (a) replacement algorithm; (b) renewal algorithm; (c) updation; (d) none of
the above.
\end{exercise}
\begin{solution}
\textbf{(a)}.
\end{solution}

\begin{exercise}[Lecturer 2]
An L2 cache helps to: (a) reduce the miss rate of L1; (b) reduce the hit time of
L1; (c) double the L1 size; (d) reduce the miss penalty of L1.
\end{exercise}
\begin{solution}
\textbf{(d)}: an L1 miss is served by L2 (a few cycles) instead of memory (about
100 cycles). L1's own miss rate and hit time do not change.
\end{solution}

\begin{exercise}[Lecturer 8--11]\label{l:cache-trace}
A CPU has an 8~KB, 2-way set-associative data cache with 16-byte lines and 32-bit
addresses.
(8) Offset bits: (a) 2; (b) 4; (c) 8; (d) 16.
(9) Index bits: (a) 2; (b) 4; (c) 8; (d) 16.
(10) Tag bits: (a) 16; (b) 18; (c) 20; (d) 22.
(11) After reset, how many misses for \texttt{0x00000000}, \texttt{0x00000014},
\texttt{0x00000008}, \texttt{0x00000010}, \texttt{0x00000020}? (a) 2; (b) 3; (c) 4;
(d) 5.
\end{exercise}
\begin{solution}
(8) \textbf{(b)}: $\log_216=4$. (9) \textbf{(c)}: $S=\dfrac{8192}{16\times2}=256$,
$\log_2256=8$. (10) \textbf{(c)}: $32-8-4=20$.

(11) \textbf{(b) 3}. Block address $=$ addr$/16$:
\begin{center}
\small
\begin{tabular}{@{}lccccc@{}}
\toprule
Address & \texttt{0x00} & \texttt{0x14} & \texttt{0x08} & \texttt{0x10} & \texttt{0x20} \\
Block (set) & 0 & 1 & 0 & 1 & 2 \\
Result & miss & miss & hit & hit & miss \\
\bottomrule
\end{tabular}
\end{center}
\texttt{0x08} is in the same 16-byte block as \texttt{0x00}, and \texttt{0x10} in the
same block as \texttt{0x14} (spatial locality). All three misses are compulsory.
\end{solution}

% ----------------------------------------------------------------------------
\subsection{Exam-style questions}
Solutions are in Section~\ref{sec:cache-sol}.

\begin{exercise}[Address breakdown]\label{q:c-addr}
A 32~KiB cache has 64-byte blocks and 32-bit byte addresses.
\begin{enumerate}[label=(\alph*),leftmargin=*]
  \item Give the tag, index and offset widths if it is direct mapped, 4-way set
  associative, and fully associative.
  \item For the 4-way write-back version, find the total storage (data, tag, valid
  and dirty bits).
  \item In the 4-way cache, which set and tag does address \texttt{0x12345678} map
  to?
\end{enumerate}
\end{exercise}

\begin{exercise}[Replacement policies]\label{q:c-repl}
One set of a 4-way set-associative cache, initially empty, receives these block
references:
\begin{center}
\texttt{0 1 2 3 0 4 1 0 5 2 3 0}
\end{center}
\begin{enumerate}[label=(\alph*),leftmargin=*]
  \item Trace LRU, FIFO, Tree-LRU and Clock. Mark hits and misses and give the
  victim of each miss. (Fill empty ways in order 0--3. Tree-LRU bits are
  $b_0b_1b_2$, initially 000; Clock gives new blocks bit 1.)
  \item Count the misses of each policy. Where does Tree-LRU differ from LRU?
  \item How many state bits per set does each policy need?
\end{enumerate}
\end{exercise}

\begin{exercise}[The 3 Cs]\label{q:c-3c}
A direct-mapped cache has 4 one-word blocks (index $=$ word address mod 4) and
starts empty. Classify every access to the word addresses
\texttt{0 4 0 1 2 3 5 0 4 1} as hit, compulsory, capacity or conflict miss.
\end{exercise}

\begin{exercise}[Write policies]\label{q:c-write}
A direct-mapped cache of 4 one-word blocks starts empty. The CPU performs
\texttt{W0 R0 W4 W4 R0 W8 R8 W0} (W $=$ write, R $=$ read, word addresses). For
(a) write-through with no-write allocate and (b) write-back with write allocate,
count the misses, the blocks read from memory and the words written to memory.
Which blocks are dirty at the end of (b)?
\end{exercise}

\begin{exercise}[Multilevel cache]\label{q:c-amat}
L1 hit time 1 cycle, L1 miss rate 2\%; L2 hit time 10 cycles, L2 local miss rate
40\%; memory 100 cycles. There are 1.3 memory accesses per instruction and
$\CPI_{exec}=1$.
\begin{enumerate}[label=(\alph*),leftmargin=*]
  \item Find the AMAT with and without L2, and the L2 global miss rate.
  \item Find the CPI with and without L2 and the speedup due to L2.
  \item Which of the 3 Cs does L2 affect? Which cache parameter does it improve?
\end{enumerate}
\end{exercise}

% ----------------------------------------------------------------------------
\subsection{Solutions}
\label{sec:cache-sol}

\begin{noot}{Exercise~\ref{q:c-addr}}
\begin{enumerate}[label=(\alph*),leftmargin=*]
  \item Offset $=\log_264=6$; $C/B=512$ blocks.
  Direct mapped: index 9, tag $32-9-6=17$.
  4-way: $S=128$, index 7, tag 19.
  Fully associative: index 0, tag 26.
  \item Per block: $512$ data $+19$ tag $+1$ valid $+1$ dirty $=533$ bits;
  $512\times533=272{,}896$ bits $=34{,}112$ bytes (about 4\% overhead).
  \item \texttt{0x12345678}$/64=$ \texttt{0x48D159}; set $=$ \texttt{0x48D159} mod
  128 $=$ \texttt{0x59} $=89$; tag $=$ \texttt{0x48D159}$/128=$ \texttt{0x91A2}.
\end{enumerate}
\end{noot}

\begin{noot}{Exercise~\ref{q:c-repl}}
(a) M $=$ miss (victim in brackets), H $=$ hit.
\begin{center}
\small
\setlength{\tabcolsep}{3pt}
\begin{tabular}{@{}lcccccccccccc@{}}
\toprule
Ref & 0 & 1 & 2 & 3 & 0 & 4 & 1 & 0 & 5 & 2 & 3 & 0 \\
\midrule
LRU & M & M & M & M & H & M(1) & M(2) & H & M(3) & M(4) & M(1) & H \\
FIFO & M & M & M & M & H & M(0) & H & M(1) & M(2) & M(3) & M(4) & H \\
Tree-LRU & M & M & M & M & H & M(2) & H & H & M(3) & M(1) & M(4) & H \\
\quad bits $b_0b_1b_2$ & 110 & 100 & 001 & 000 & 110 & 011 & 101 & 111 & 010 & 100 & 001 & 111 \\
Clock & M & M & M & M & H & M(0) & H & M(2) & M(3) & M(1) & M(4) & H \\
\bottomrule
\end{tabular}
\end{center}
Tree-LRU detail at reference 4: the bits are 110 ($b_0=1$: victim in pair
$\{2,3\}$; $b_2=0$: way 2), so block 2 is evicted, while true LRU evicts 1. Clock at
reference 4: all bits are 1, so the pointer clears all four and comes back to way 0,
evicting 0.

(b) LRU 9, FIFO 9, Tree-LRU 8, Clock 9. Tree-LRU differs at reference 4 (it keeps 1,
which is reused next) and from then on holds different blocks. Here the
approximation happens to do better; in general it is slightly worse than LRU.

(c) Exact LRU: $\lceil\log_2 4!\rceil=5$ bits (often 8: a 2-bit age per way).
FIFO: a 2-bit pointer. Tree-LRU: 3 bits. Clock: 4 reference bits $+$ a 2-bit
pointer. Random: none.
\end{noot}

\begin{noot}{Exercise~\ref{q:c-3c}}
\begin{center}
\small
\begin{tabular}{@{}lcccccccccc@{}}
\toprule
Address & 0 & 4 & 0 & 1 & 2 & 3 & 5 & 0 & 4 & 1 \\
Index & 0 & 0 & 0 & 1 & 2 & 3 & 1 & 0 & 0 & 1 \\
Result & comp. & comp. & confl. & comp. & comp. & comp. & comp. & hit & cap. & cap. \\
\bottomrule
\end{tabular}
\end{center}
The third access (0) misses only because 4 took index 0; a 4-block fully
associative LRU cache would still hold 0, so it is a conflict miss. The last two
(4 and 1) also miss in the fully associative cache, which has meanwhile evicted them
(too many distinct blocks since their last use), so they are capacity misses.
\end{noot}

\begin{noot}{Exercise~\ref{q:c-write}}
\begin{enumerate}[label=(\alph*),leftmargin=*]
  \item Write-through, no-write allocate: writes never fill the cache. W0 miss, R0
  miss (fetch 0), W4 miss, W4 miss, R0 hit, W8 miss, R8 miss (fetch 8, replaces 0),
  W0 miss. \textbf{7 misses, 2 blocks read, 5 words written} (every write goes to
  memory).
  \item Write-back, write allocate: W0 miss (fetch 0, dirty), R0 hit, W4 miss
  (write back 0, fetch 4, dirty), W4 hit, R0 miss (write back 4, fetch 0, clean), W8
  miss (0 is clean: no write-back; fetch 8, dirty), R8 hit, W0 miss (write back 8,
  fetch 0, dirty). \textbf{5 misses, 5 blocks read, 3 words written}; block 0 is
  dirty at the end (one more write-back when it is evicted).
\end{enumerate}
\end{noot}

\begin{noot}{Exercise~\ref{q:c-amat}}
\begin{enumerate}[label=(\alph*),leftmargin=*]
  \item With L2: $1+0.02\times(10+0.4\times100)=1+0.02\times50=2$ cycles. Without:
  $1+0.02\times100=3$ cycles. Global L2 miss rate $=0.02\times0.4=0.8\%$.
  \item With L2: stalls $=1.3\times0.02\times50=1.3$, $\CPI=2.3$. Without:
  $1.3\times0.02\times100=2.6$, $\CPI=3.6$. Speedup $=3.6/2.3=1.57$.
  \item L2 does not change L1's misses (any C); it reduces the \emph{miss penalty}
  of L1, because most L1 misses now hit in L2.
\end{enumerate}
\end{noot}


% ============================================================================
\section{Virtual Memory}
\label{sec:vm}

Address translation, page tables and the TLB, paging and segmentation, page faults
and page replacement.

% ----------------------------------------------------------------------------
\subsection{What to know}

\subsubsection{Why virtual memory}
\begin{itemize}[leftmargin=*]
  \item Each process sees its own large \textbf{virtual} (logical) address space; the
  OS and hardware map it onto physical memory, with disk as backing store.
  \item Benefits: a program can be \textbf{larger than physical memory}; processes
  are \textbf{isolated} (the same virtual address in two programs maps to different
  physical addresses) and protected; pages can be \textbf{shared}; programs are
  \textbf{relocatable}; consecutive virtual pages (e.g.\ \texttt{A[n]} and
  \texttt{A[n+1]} across a page boundary) need not be adjacent in physical memory.
  \item Terms: VA, PA; VPN (virtual page number), PFN/PPN (physical frame number),
  PTE (page table entry), MMU, TLB, ASID.
\end{itemize}

\subsubsection{Paging and address translation}
Page size $P$ (e.g.\ 4~KiB $=2^{12}$, so the low 12 bits, 3 hex digits, are the
offset):
\[
  \text{VPN}=\lfloor\text{VA}/P\rfloor,\quad \text{offset}=\text{VA}\bmod P,\quad
  \text{PA}=\text{PFN}\times P+\text{offset}.
\]
\begin{itemize}[leftmargin=*]
  \item The \textbf{page table} (one per process, indexed by VPN) stores the
  \textbf{physical page number} plus valid/present, dirty, reference and protection
  bits (R/W/X, user/kernel). A page-table base register points to it.
  \item Flat table size $=2^{\text{VPN bits}}\times$ PTE size, \emph{per process}
  (32-bit VA, 4~KiB pages, 4~B PTE: $2^{20}\times4=4$~MiB).
  \item \textbf{Multilevel} table: split the VPN; only the parts in use are
  allocated. A level fits in one page when it has $\log_2(P/\text{PTE size})$ bits
  (4~KiB$/$4~B: $10+10+12$). A walk costs one memory access per level.
  \item \textbf{Inverted} table: one entry per physical frame, searched by hashing
  the VPN; size proportional to physical memory.
  \item Placement is fully associative (any page in any frame); replacement is done
  by the OS; writes are always \textbf{write-back} (dirty bit), because disk is slow.
\end{itemize}

\subsubsection{Page faults}
\begin{itemize}[leftmargin=*]
  \item A \textbf{page fault} is the exception raised when the PTE says the page is
  not in memory (a ``page-table miss''). The OS: checks the access is legal (else
  kills the process), picks a victim frame (writing it back if dirty), reads the page
  from disk, updates the PTE, and restarts the faulting instruction.
  \item Effective access time with fault rate $p$ and fault service time $s$:
  $\text{EAT}=(1-p)\,m+p\,s$. With $m=100$~ns and $s=10$~ms, $p=0.1\%$ already
  gives $\approx10{,}100$~ns, about $101\times$ slower.
  \item \textbf{Thrashing}: the working sets of the running processes do not fit in
  memory, so the system spends its time paging.
\end{itemize}

\subsubsection{TLB}
\begin{itemize}[leftmargin=*]
  \item The \textbf{TLB} is a small, fast (usually fully associative) cache of
  recent VPN $\to$ PFN translations. It speeds up \emph{translation}: without it,
  every access needs a page-table read plus the data read.
  \item TLB hit: PA $=$ PFN $\Vert$ offset at once. TLB miss: walk the page table
  (hardware or OS); if the page is present, load the PTE into the TLB and retry; if
  not, page fault. So a TLB miss \emph{may or may not} cause a page fault, but a TLB
  hit never does.
  \item $\text{EAT}=h\,(t+m)+(1-h)\,(t+k\,m+m)$ for TLB time $t$, memory time $m$,
  hit rate $h$, $k$-level page table ($k=1$: $t+2m$ on a miss).
  \item TLB reach $=$ entries $\times$ page size. On a context switch, flush the TLB
  or tag entries with an ASID.
\end{itemize}

\subsubsection{Page replacement}
\begin{itemize}[leftmargin=*]
  \item \textbf{OPT}: evict the page used furthest in the future (unrealizable;
  the lower bound). \textbf{FIFO}: evict the oldest loaded page; simple, but can
  evict a busy page and shows \textbf{Belady's anomaly} (more frames, more faults:
  \texttt{1 2 3 4 1 2 5 1 2 3 4 5} gives 9 faults with 3 frames, 10 with 4).
  \textbf{LRU}: evict the page unused for the longest time; uses temporal locality,
  costly to track exactly, never shows Belady's anomaly. \textbf{Clock} (second
  chance): approximates LRU with one reference bit per frame. LFU: evict the least
  frequently used.
  \item With $k$ empty frames, the first $k$ distinct pages always fault.
\end{itemize}

\subsubsection{Segmentation}
\begin{itemize}[leftmargin=*]
  \item Memory is divided into \textbf{variable-size} logical segments (code, data,
  heap, stack), visible to the programmer. Address $=$ (segment, offset). Each
  segment-table entry has \textbf{base}, \textbf{limit} and protection:
  if offset $<$ limit then PA $=$ base $+$ offset, else trap.
  \item Paging vs.\ segmentation: fixed vs.\ variable size; internal fragmentation
  (half a page per region on average) vs.\ \textbf{external fragmentation} (holes,
  needs compaction; placement first/best/worst fit); per-page vs.\ per-region
  protection and sharing.
  \item \textbf{Segmented paging}: each segment has its own page table (Multics,
  x86 protected mode). x86 real mode: PA $=$ segment $\times16+$ offset.
\end{itemize}

\begin{center}
\small
\begin{tabular}{@{}lll@{}}
\toprule
& L1 cache & Virtual memory \\
\midrule
Unit & block, 16--128 B & page, 4--64 KiB \\
Hit time & 1--3 cycles & 100--200 cycles \\
Miss penalty & 8--200 cycles & $10^6$--$10^7$ cycles \\
Miss rate & 0.1--10\% & 0.00001--0.001\% \\
Placement & direct / set associative & fully associative \\
Replacement & hardware & OS (LRU approximation) \\
Writes & write-through or write-back & always write-back \\
\bottomrule
\end{tabular}
\end{center}

% ----------------------------------------------------------------------------
\subsection{Common traps}
\begin{enumerate}[leftmargin=*]
  \item TLB miss $\neq$ page fault. The page table stores \emph{physical} page
  numbers.
  \item Convert ms to ns ($\times10^6$) for page-fault service times.
  \item Page-table size is per process: multiply by the number of processes if asked.
  \item In a $k$-level table, a TLB miss costs $k$ extra memory accesses.
  \item FIFO: hits do not change the order. LRU: every access moves the page to most
  recent.
  \item Segmentation: check the limit \emph{before} adding the base.
\end{enumerate}

% ----------------------------------------------------------------------------
\subsection{Lecturer's review questions}

\begin{exercise}[Lecturer 1--4]
(1) In virtual memory: (a) a program may be larger than physical memory; (b) the same
address in 2 different programs can be different; (c) the elements \texttt{A[n]} and
\texttt{A[n+1]} need not lie side by side in physical memory; (d) all of the above.

(2) The binary addresses issued for data or instructions are called: (a) physical
addresses; (b) locations; (c) relocatable addresses; (d) logical addresses.

(3) The page table stores: (a) data of the physical page; (b) virtual page data;
(c) physical page addresses; (d) virtual page addresses.

(4) Reading data at an address in a paging system takes: (a) 1 physical address
lookup and 1 physical data read; (b) 1 virtual address lookup and 1 virtual data
read; (c) 1 virtual address lookup and 1 physical data read; (d) 1 physical data
read, and if it misses, 1 more virtual data read.
\end{exercise}
\begin{solution}
(1) \textbf{(d)}: all three are what virtual memory gives (size, isolation,
non-contiguous mapping).
(2) \textbf{(d)}: the CPU issues logical (virtual) addresses; the MMU translates
them.
(3) \textbf{(c)}: each PTE holds the physical page (frame) number of a virtual page,
plus status bits.
(4) \textbf{(c)}: look up the virtual address (VPN) in the TLB or page table, then
read the data at the physical address.
\end{solution}

\begin{exercise}[Lecturer 5--9]
(5) A page-fault exception: (a) indicates a cache miss; (b) indicates a page-table
miss; (c) means the page size is larger than standard; (d) means physical and
virtual pages have different sizes.

(6) The TLB is used to: (a) accelerate physical data access; (b) accelerate hard
disk access; (c) accelerate the translation from virtual to physical addresses;
(d) all of the above.

(7) When a TLB miss occurs: (a) a page fault occurs; (b) no page fault occurs;
(c) the system is reset; (d) a page fault may or may not occur.

(8) In FIFO page replacement, the page replaced is: (a) the newest; (b) the oldest;
(c) a random one; (d) none of the mentioned.

(9) The algorithm that replaces the block not referenced for the longest while is:
(a) LRU; (b) ORF; (c) direct; (d) both LRU and ORF.
\end{exercise}
\begin{solution}
(5) \textbf{(b)}: the PTE shows the page is not in memory.
(6) \textbf{(c)}.
(7) \textbf{(d)}: the page-table walk may find the page present (just refill the TLB)
or not present (page fault).
(8) \textbf{(b)}.
(9) \textbf{(a)} least recently used.
\end{solution}

\begin{exercise}[Lecturer 10--11]\label{l:vm-repl}
A system has only three physical pages, initially unused. For the virtual page
references
\begin{center}
\texttt{1 2 1 1 3 2 1 4 3 1 1 2 4 1 5 6 2 1}
\end{center}
find the number of page faults with (10) FIFO and (11) LRU. Choices: (a) 9; (b) 10;
(c) 11; (d) 12.
\end{exercise}
\begin{solution}
F $=$ fault (evicted page in brackets), $\cdot$ $=$ hit.
\begin{center}
\small
\setlength{\tabcolsep}{2.5pt}
\begin{tabular}{@{}l*{18}{c}@{}}
\toprule
Ref & 1 & 2 & 1 & 1 & 3 & 2 & 1 & 4 & 3 & 1 & 1 & 2 & 4 & 1 & 5 & 6 & 2 & 1 \\
\midrule
FIFO & F & F & $\cdot$ & $\cdot$ & F & $\cdot$ & $\cdot$ & F(1) & $\cdot$ & F(2) & $\cdot$ & F(3) & $\cdot$ & $\cdot$ & F(4) & F(1) & $\cdot$ & F(2) \\
LRU & F & F & $\cdot$ & $\cdot$ & F & $\cdot$ & $\cdot$ & F(3) & F(2) & $\cdot$ & $\cdot$ & F(4) & F(3) & $\cdot$ & F(2) & F(4) & F(1) & F(5) \\
\bottomrule
\end{tabular}
\end{center}
(10) FIFO: \textbf{9 faults, (a)}. (11) LRU: \textbf{11 faults, (c)}.

Why LRU loses here: at reference 4 (frames 1, 2, 3, last used in the order 3, 2, 1)
LRU evicts 3, which is needed right away; FIFO evicts 1, the oldest load, and 1 is
not needed until two references later. LRU is usually better, but not on every
string. (OPT gives 7 faults; see Exercise~\ref{q:vm-repl}.)
\end{solution}

% ----------------------------------------------------------------------------
\subsection{Exam-style questions}
Solutions are in Section~\ref{sec:vm-sol}.

\begin{exercise}[Address translation]\label{q:vm-trans}
32-bit virtual addresses, 4~KiB pages. Part of the page table (hex):
\begin{center}
\small
\begin{tabular}{@{}lll@{}}
\toprule
VPN & Valid & PFN \\
\midrule
\texttt{00400} & 1 & \texttt{0A1} \\
\texttt{00401} & 0 & (on disk) \\
\texttt{00403} & 1 & \texttt{1F2} \\
\bottomrule
\end{tabular}
\end{center}
Translate \texttt{0x00403A7C}, \texttt{0x00400008} and \texttt{0x00401010}.
\end{exercise}

\begin{exercise}[Page table size]\label{q:vm-size}
\leavevmode
\begin{enumerate}[label=(\alph*),leftmargin=*]
  \item 32-bit VA, 4~KiB pages, 4-byte PTEs. Find the flat page-table size per
  process and for 100 processes.
  \item With a two-level table ($10+10+12$), a process uses 8~MiB of code and data
  at the bottom of its address space and 1~MiB of stack at the top. How much memory
  do its page tables take?
  \item With 48-bit VA, 4~KiB pages and 8-byte PTEs, why is a flat table impossible?
  How many levels are needed if each level must fit in one page?
\end{enumerate}
\end{exercise}

\begin{exercise}[TLB and page faults]\label{q:vm-eat}
TLB lookup 10~ns, memory access 100~ns, TLB hit rate 95\%, two-level page table.
\begin{enumerate}[label=(\alph*),leftmargin=*]
  \item Find the effective access time (no page faults).
  \item Page faults happen on $10^{-6}$ of accesses and take 8~ms. Find the new
  effective access time.
  \item Which combinations are possible: (i) TLB hit, page fault; (ii) TLB miss,
  page present, data in cache; (iii) TLB miss, page fault, data in cache;
  (iv) TLB miss, page present, data not in cache?
\end{enumerate}
\end{exercise}

\begin{exercise}[More replacement]\label{q:vm-repl}
For the reference string of Lecturer question 10 with 3 frames:
\begin{enumerate}[label=(\alph*),leftmargin=*]
  \item Find the faults with OPT.
  \item Find the faults with Clock (new pages get reference bit 1; the pointer
  starts at frame 0 and moves past each newly loaded page).
  \item What is Belady's anomaly, and which of FIFO, LRU, OPT, Clock can show it?
\end{enumerate}
\end{exercise}

\begin{exercise}[Segmentation]\label{q:vm-seg}
Segment table: segment 0 base 1400, limit 1000; segment 1 base 6300, limit 400;
segment 2 base 4300, limit 400. Translate (2, 53), (1, 500), (0, 999).
Give one advantage and one drawback of segmentation compared with paging.
\end{exercise}

% ----------------------------------------------------------------------------
\subsection{Solutions}
\label{sec:vm-sol}

\begin{noot}{Exercise~\ref{q:vm-trans}}
Offset $=$ low 12 bits (3 hex digits), VPN $=$ upper 20 bits.
\texttt{0x00403A7C}: VPN \texttt{00403}, offset \texttt{A7C} $\to$ PA
\texttt{0x1F2A7C}.
\texttt{0x00400008}: VPN \texttt{00400}, offset \texttt{008} $\to$ PA
\texttt{0x0A1008}.
\texttt{0x00401010}: VPN \texttt{00401} is not valid $\to$ \textbf{page fault}; the
OS loads the page, updates the PTE and restarts the access.
\end{noot}

\begin{noot}{Exercise~\ref{q:vm-size}}
\begin{enumerate}[label=(\alph*),leftmargin=*]
  \item $2^{20}$ PTEs $\times4$~B $=4$~MiB per process; 400~MiB for 100 processes.
  \item Each second-level table maps $2^{10}\times4$~KiB $=4$~MiB. Code and data need
  2 tables, the stack 1, plus the root: $4\times4$~KiB $=16$~KiB (instead of 4~MiB).
  \item VPN $=48-12=36$ bits: $2^{36}\times8$~B $=512$~GiB per process. A page holds
  $4096/8=512=2^9$ PTEs, so $\lceil36/9\rceil=4$ levels (as in x86-64).
\end{enumerate}
\end{noot}

\begin{noot}{Exercise~\ref{q:vm-eat}}
\begin{enumerate}[label=(\alph*),leftmargin=*]
  \item Hit: $10+100=110$~ns. Miss: $10+2\times100+100=310$~ns.
  EAT $=0.95\times110+0.05\times310=104.5+15.5=120$~ns.
  \item $+10^{-6}\times8\times10^6$~ns $=+8$~ns: about 128~ns.
  \item (i) Impossible: a translation is in the TLB only if the page is in memory.
  (ii) Possible. (iii) Impossible: data cannot be cached if its page is not in
  memory. (iv) Possible.
\end{enumerate}
\end{noot}

\begin{noot}{Exercise~\ref{q:vm-repl}}
\begin{enumerate}[label=(\alph*),leftmargin=*]
  \item OPT: faults at 1, 2, 3, 4 (evict 2, next used furthest), 2 (evict 3), 5
  (evict 4, never used again), 6 (evict 5): \textbf{7 faults}.
  \item Clock: faults at 1, 2, 3, 4 (all bits 1: clear them all, evict 1), 1 (evict
  2), 2 (evict 3), 5 (evict 4), 6 (evict 1), 1 (evict 2): \textbf{9 faults}.
  \item More frames giving \emph{more} faults. Only FIFO (and Clock, which degrades
  to FIFO when all bits are set) can show it; LRU and OPT are stack algorithms and
  never do.
\end{enumerate}
\end{noot}

\begin{noot}{Exercise~\ref{q:vm-seg}}
(2, 53): $53<400$, PA $=4300+53=4353$. (1, 500): $500\ge400$ $\to$ trap, no address
is formed. (0, 999): $999<1000$, PA $=1400+999=2399$. Advantage: segments match
logical regions, so protection and sharing are natural, with no internal
fragmentation. Drawback: variable sizes cause external fragmentation (needing
compaction), and growing a segment is hard.
\end{noot}


% ============================================================================
\section{Multiprocessors and Instruction-Level Parallelism}
\label{sec:mp}

Flynn's classification, UMA and NUMA, superpipelining, superscalar and VLIW
processors, dependences and issue/completion order, multithreading, cache coherence
(snoopy and directory protocols, MSI, MESI) and synchronization (spin locks,
barriers).

\begin{remark}
The course slides for this part are not in this repository; this section follows
H\&P (Ch.~3 and~5) and the topic list in the lecturer's review slides. Check terms
against the lecture slides.
\end{remark}

% ----------------------------------------------------------------------------
\subsection{What to know}

\subsubsection{Classifying parallel machines}
\begin{center}
\small
\begin{tabular}{@{}lp{11.2cm}@{}}
\toprule
\multicolumn{2}{@{}l}{\textbf{Flynn's classification} (instruction streams $\times$ data streams)} \\
\midrule
SISD & one instruction stream, one data stream: a uniprocessor (may still be pipelined or superscalar). \\
SIMD & one instruction applied to many data: vector processors, multimedia extensions (SSE, AVX), GPUs. \\
MISD & several instruction streams on one data stream: \textbf{no commercial machine}; not used in practice. \\
MIMD & independent processors on their own data: multicores, multiprocessors, clusters. The most general and common. \\
\bottomrule
\end{tabular}
\end{center}
\begin{itemize}[leftmargin=*]
  \item \textbf{UMA} (uniform memory access), also \textbf{SMP} or centralized
  shared memory: all processors share one memory with the same access time; simple,
  but the bus/memory limits scaling (a few dozen cores).
  \item \textbf{NUMA} (non-uniform memory access), also distributed shared memory
  (DSM): memory is distributed among the nodes but still one address space; local
  accesses are faster than remote ones; scales further. Average access time
  $=f_{local}\,t_{local}+(1-f_{local})\,t_{remote}$.
  \item \textbf{Message passing} (clusters, multicomputers): private memories,
  explicit send/receive.
\end{itemize}

\subsubsection{Instruction-level parallelism (ILP)}
\begin{itemize}[leftmargin=*]
  \item \textbf{Superpipelining}: split stages into more, shorter stages, giving a
  higher clock rate. Costs: larger branch and load-use penalties (in cycles), more
  latch overhead.
  \item \textbf{Superscalar}: issue several instructions per cycle (ideal
  $\CPI=1/n$, $\text{IPC}=n$); hardware checks dependences dynamically, statically
  scheduled (in order) or dynamically scheduled (out of order); binary compatible.
  \item \textbf{VLIW}: the \emph{compiler} packs independent operations into one
  long instruction with a fixed slot per functional unit; no dynamic hazard
  detection, so simpler hardware. Drawbacks: \texttt{nop}s in unused slots increase
  \textbf{code size}; latencies are built into the code, so binaries are not
  compatible across generations; cache misses stall the whole bundle. (EPIC/IA-64 is
  a VLIW variant.)
\end{itemize}

\subsubsection{Dependences and issue order}
\begin{itemize}[leftmargin=*]
  \item \textbf{Data (true) dependence}: $j$ uses a result of $i$ (RAW hazard).
  \textbf{Name dependences}: \emph{anti-dependence} ($j$ writes a register $i$ reads:
  WAR) and \emph{output dependence} (both write it: WAW); no value flows, only the
  register name is shared. \textbf{Control dependence}: on a branch outcome.
  \item \textbf{Register renaming} (Tomasulo's reservation stations, or physical
  register files) gives each new value a fresh register and \textbf{removes WAR and
  WAW}; RAW must still be respected.
\end{itemize}
\begin{center}
\small
\begin{tabular}{@{}p{4.4cm}p{9.8cm}@{}}
\toprule
Policy & Hazards to check \\
\midrule
In-order issue, in-order completion & RAW only (like the 5-stage pipeline); a slow instruction blocks all behind it. \\
In-order issue, out-of-order completion & RAW and \textbf{WAW} (a later, shorter instruction can write first). Precise exceptions become hard. \\
Out-of-order issue, out-of-order completion & RAW, WAW and \textbf{WAR} (a later instruction may write before an earlier one reads). Needs an instruction window, renaming, and a reorder buffer for in-order \emph{commit} (precise exceptions, speculation). \\
\bottomrule
\end{tabular}
\end{center}
Scoreboarding issues out of order and stalls on WAR/WAW; Tomasulo's algorithm renames
registers to remove them; hardware speculation adds a reorder buffer.

\subsubsection{Multithreading}
Several threads share one core; each has its own PC and registers.
\begin{itemize}[leftmargin=*]
  \item \textbf{Fine-grained}: switch threads every cycle (round robin), skipping
  stalled threads; hides short and long stalls but slows each single thread.
  \item \textbf{Coarse-grained}: switch only on long stalls (L2 miss); little impact
  on one thread, but the pipeline must refill on a switch, so short stalls are not
  hidden.
  \item \textbf{Simultaneous multithreading (SMT)}: on a superscalar out-of-order
  core, issue instructions \emph{from several threads in the same cycle}, exploiting
  \textbf{both thread-level and instruction-level parallelism} (Intel
  Hyper-Threading).
\end{itemize}

\subsubsection{Cache coherence}
\begin{itemize}[leftmargin=*]
  \item With private caches, two processors can hold different values for one
  address. A system is \textbf{coherent} if a read returns the last value written
  (writes to one location are seen in the same order by all processors).
  Consistency is about the order of accesses to \emph{different} locations.
  \item \textbf{Snoopy protocols}: every cache watches (snoops) a shared bus;
  broadcast; good for small UMA machines. \textbf{Write-invalidate} (the usual):
  a writer first invalidates the other copies. \textbf{Write-update}: broadcast the new
  value.
  \item \textbf{Directory protocols}: a directory (at each block's home node) records
  the block's state and which caches share it (a bit vector); requests go point to
  point to the home, which sends invalidations only to the sharers. Scales to NUMA.
  Directory states: uncached, shared, exclusive/modified.
  \item Coherence misses (the fourth C): \emph{true sharing} (the word is really
  communicated) and \emph{false sharing} (different words in the same block).
\end{itemize}

\begin{center}
\small
\begin{tabular}{@{}p{2.6cm}p{11.6cm}@{}}
\toprule
State & Meaning \\
\midrule
M (Modified) & the only copy, dirty (memory is stale); must supply data and write back. \\
E (Exclusive) & MESI only: the only copy, clean (same as memory). \\
S (Shared) & clean, possibly in other caches. \\
I (Invalid) & not present / not valid. \\
\bottomrule
\end{tabular}
\end{center}
\begin{center}
\small
\begin{tabular}{@{}lll@{}}
\toprule
Event & MSI & MESI \\
\midrule
Read miss & BusRd $\to$ S & BusRd $\to$ E if no other copy, else S \\
Write miss & BusRdX $\to$ M & BusRdX $\to$ M \\
Write hit in S & BusUpgr (invalidate) $\to$ M & BusUpgr $\to$ M \\
Write hit in E & -- & \textbf{silent} E $\to$ M (no bus transaction) \\
Snoop BusRd, in M & flush (write back, supply data) $\to$ S & flush $\to$ S \\
Snoop BusRd, in E & -- & $\to$ S \\
Snoop BusRdX/BusUpgr & M: flush $\to$ I; S $\to$ I & M: flush $\to$ I; E, S $\to$ I \\
Evict & M: write back; S: silent & M: write back; E, S: silent \\
\bottomrule
\end{tabular}
\end{center}
The \textbf{E state} lets a processor that read a private line write it later without
a bus transaction: the common case for single-threaded code and private data.

\subsubsection{Synchronization}
\begin{itemize}[leftmargin=*]
  \item Built on an \textbf{atomic read-modify-write}: test-and-set, atomic exchange,
  fetch-and-increment, or the \texttt{ll}/\texttt{sc} pair (Section~\ref{sec:isa}).
  \item \textbf{Spin lock}: a processor loops until it acquires the lock. Spin on a
  \emph{cached copy} (plain loads, ``test-and-test-and-set'') and attempt the atomic
  operation only when the lock looks free; this keeps the spinning off the bus. With
  \texttt{ll}/\texttt{sc}:
\begin{lstlisting}
lockit: ll   $t0, 0($s0)       # read lock
        bne  $t0, $zero, lockit # spin while it is held (1)
        addi $t0, $zero, 1
        sc   $t0, 0($s0)       # try to set it to 1
        beq  $t0, $zero, lockit # sc failed: retry
        ...                     # critical section
        sw   $zero, 0($s0)     # release
\end{lstlisting}
  \item \textbf{Barrier}: no thread continues until all $n$ have arrived. Simple
  version: a lock-protected counter; the last arrival resets the counter and sets
  \texttt{release}; the others spin on \texttt{release}. Reusing it in a loop can
  deadlock (a fast thread re-enters before the slow ones have seen
  \texttt{release}); the fix is a \textbf{sense-reversing barrier}: each thread
  flips a private \texttt{local\_sense}, and the last arrival sets
  \texttt{release = local\_sense}.
\end{itemize}

% ----------------------------------------------------------------------------
\subsection{Common traps}
\begin{enumerate}[leftmargin=*]
  \item MISD is the class with no practical machines; SIMD is widely used (vectors,
  GPUs).
  \item Renaming removes WAR and WAW, never RAW.
  \item WAW appears with out-of-order \emph{completion}; WAR needs out-of-order
  \emph{issue} (or reads after later writes).
  \item VLIW: the compiler, not the hardware, finds the parallelism.
  \item SMT issues from several threads in the \emph{same} cycle; fine-grained
  switches every cycle but issues from one thread per cycle.
  \item MESI E $\to$ M needs no bus transaction; S $\to$ M does.
  \item A read snoop on an M line forces a write-back (flush).
\end{enumerate}

% ----------------------------------------------------------------------------
\subsection{Lecturer's review questions}

\begin{exercise}[Lecturer 1]
In the MESI protocol, what is the principal purpose of the E state?
(a) to avoid a second bus transaction when a single-threaded program writes a line
it previously read; (b) to identify the core responsible for responding to BusRd
requests from other cores; (c) to avoid the need for write-back when evicting an
exclusive line whose value is still correct in memory; (d) to help identify the
oldest line in a set when a victim must be chosen.
\end{exercise}
\begin{solution}
\textbf{(a)}. A read with no other sharers loads the line in E; a later write changes
E $\to$ M silently. With MSI the line would be in S, and the write would need a
BusUpgr. (c) is also true of S lines, so it is not what E adds.
\end{solution}

\begin{exercise}[Lecturer 2]
Register renaming helps to eliminate: (a) control hazards; (b) structural hazards;
(c) RAW and WAR data hazards; (d) WAR and WAW data hazards.
\end{exercise}
\begin{solution}
\textbf{(d)}: renaming removes name dependences. RAW is a true dependence; the
value must still be produced before it is used.
\end{solution}

\begin{exercise}[Lecturer 4]
The main feature of simultaneous multithreading is that it exploits both:
(a) thread-level and instruction-level parallelism; (b) coarse-grained multithreading
and speculation; (c) fine-grained multithreading and multilevel cache hierarchies;
(d) all of the above.
\end{exercise}
\begin{solution}
\textbf{(a)}: several threads (TLP) issue in the same cycle to a wide out-of-order
core (ILP).
\end{solution}

\begin{exercise}[Lecturer 5]
In Flynn's classification, which model is not much used in practice?
(a) SISD; (b) SIMD; (c) MIMD; (d) MISD.
\end{exercise}
\begin{solution}
\textbf{(d)} MISD.
\end{solution}

\begin{exercise}[Lecturer 6]
A VLIW machine: (a) relies on the compiler to schedule all instructions;
(b) dynamically detects hazards to reduce their delay penalty; (c) improves binary
compatibility between generations of the architecture; (d) improves code size.
\end{exercise}
\begin{solution}
\textbf{(a)}. (b) describes superscalar out-of-order hardware; (c) and (d) are VLIW's
weaknesses (latencies fixed in the code, \texttt{nop}-filled slots).
\end{solution}

% ----------------------------------------------------------------------------
\subsection{Exam-style questions}
Solutions are in Section~\ref{sec:mp-sol}.

\begin{exercise}[MSI vs.\ MESI trace]\label{q:mp-mesi}
Three processors with private write-back caches share a snoopy bus. Block $X$ is
initially in no cache. For each access give the state of $X$ in P1, P2, P3 and the bus
action, under MSI and under MESI:
\begin{center}
P1 reads $X$; P2 reads $X$; P1 writes $X$; P3 reads $X$; P2 writes $X$; P1 reads $X$.
\end{center}
Then P1 reads block $Y$ (cached nowhere) and writes it. Count the bus transactions for
$Y$ under each protocol.
\end{exercise}

\begin{exercise}[Dependences and issue policies]\label{q:mp-dep}
\leavevmode
\begin{lstlisting}
I1: div  r1, r2, r3     # long latency
I2: add  r4, r1, r5
I3: sub  r5, r6, r7
I4: mul  r4, r6, r8
I5: and  r9, r5, r4
\end{lstlisting}
\begin{enumerate}[label=(\alph*),leftmargin=*]
  \item List all true, anti and output dependences.
  \item With in-order issue and in-order completion, what happens while I1 executes?
  \item With in-order issue and out-of-order completion, which new hazard must be
  checked? Give the instruction pair.
  \item With out-of-order issue, which hazards appear if I3 or I4 runs before I2?
  \item Rename the registers so only true dependences remain.
\end{enumerate}
\end{exercise}

\begin{exercise}[VLIW scheduling]\label{q:mp-vliw}
A 2-slot VLIW has one memory slot and one ALU slot per instruction word. A load's
result can be used 2 cycles after it issues; an ALU result 1 cycle after. Registers
are read before they are written within one word. Schedule
\begin{lstlisting}
lw   r1, 0(r10)
lw   r2, 4(r10)
add  r3, r1, r2
sw   r3, 8(r10)
addi r10, r10, 12
\end{lstlisting}
in as few words as possible. How many words and \texttt{nop}s are there? What would a
superscalar processor do differently?
\end{exercise}

\begin{exercise}[Spin locks and barriers]\label{q:mp-sync}
\leavevmode
\begin{enumerate}[label=(\alph*),leftmargin=*]
  \item In the \texttt{ll}/\texttt{sc} spin lock above, why spin on the value with
  \texttt{bne} (line (1)) before trying \texttt{sc}? What would happen on the bus if
  every processor executed an atomic exchange in a tight loop?
  \item Explain why the simple counter barrier can fail when used twice in a loop,
  and how sense reversal fixes it.
\end{enumerate}
\end{exercise}

\begin{exercise}[UMA, NUMA and multithreading]\label{q:mp-misc}
\leavevmode
\begin{enumerate}[label=(\alph*),leftmargin=*]
  \item On a NUMA machine a local access takes 20~ns and a remote one 100~ns. Find
  the average access time if 80\% of accesses are local, and if only 50\% are.
  \item A core runs two threads; thread A stalls on long cache misses, thread B has
  many short 1-cycle dependence stalls. Which multithreading style hides each kind of
  stall?
  \item Why does superpipelining raise the branch penalty in cycles?
\end{enumerate}
\end{exercise}

% ----------------------------------------------------------------------------
\subsection{Solutions}
\label{sec:mp-sol}

\begin{noot}{Exercise~\ref{q:mp-mesi}}
States are listed as P1 P2 P3.
\begin{center}
\small
\begin{tabular}{@{}lllll@{}}
\toprule
Access & MSI states & MSI bus & MESI states & MESI bus \\
\midrule
P1 rd $X$ & S I I & BusRd & \textbf{E} I I & BusRd \\
P2 rd $X$ & S S I & BusRd & S S I & BusRd (P1 E $\to$ S) \\
P1 wr $X$ & M I I & BusUpgr & M I I & BusUpgr \\
P3 rd $X$ & S I S & BusRd, P1 flushes & S I S & BusRd, P1 flushes \\
P2 wr $X$ & I M I & BusRdX & I M I & BusRdX \\
P1 rd $X$ & S S I & BusRd, P2 flushes & S S I & BusRd, P2 flushes \\
\bottomrule
\end{tabular}
\end{center}
Both use 6 bus transactions for $X$ (with 2 write-backs). For $Y$: MSI uses BusRd
(S) then BusUpgr (M), 2 transactions; MESI uses BusRd (E) then a silent E $\to$ M,
1 transaction. This is the purpose of E.
\end{noot}

\begin{noot}{Exercise~\ref{q:mp-dep}}
\begin{enumerate}[label=(\alph*),leftmargin=*]
  \item True (RAW): I1$\to$I2 (r1), I3$\to$I5 (r5), I4$\to$I5 (r4). Anti (WAR):
  I2$\to$I3 (r5: I2 reads, I3 writes). Output (WAW): I2$\to$I4 (r4).
  \item I2 waits for r1, and because issue is in order, I3--I5 wait behind it even
  though I3 and I4 are independent of I1.
  \item WAW between I2 and I4: if I4 (\texttt{mul}) completes before I2, I2's later
  write would leave the wrong value in r4 for I5.
  \item I3 before I2: WAR, I3 overwrites r5 before I2 reads it. I4 before I2: WAW on
  r4.
  \item \texttt{I3: sub p5,r6,r7}; \texttt{I4: mul p4,r6,r8};
  \texttt{I5: and r9,p5,p4}. Now I3 and I4 can run at once during the \texttt{div};
  only I1$\to$I2 and I3, I4$\to$I5 remain.
\end{enumerate}
\end{noot}

\begin{noot}{Exercise~\ref{q:mp-vliw}}
\begin{center}
\small
\begin{tabular}{@{}lll@{}}
\toprule
Word & Memory slot & ALU slot \\
\midrule
1 & \texttt{lw r1,0(r10)} & \texttt{nop} \\
2 & \texttt{lw r2,4(r10)} & \texttt{nop} \\
3 & \texttt{nop} & \texttt{nop} \\
4 & \texttt{nop} & \texttt{add r3,r1,r2} \\
5 & \texttt{sw r3,8(r10)} & \texttt{addi r10,r10,12} \\
\bottomrule
\end{tabular}
\end{center}
5 words, 5 \texttt{nop}s (50\% of the slots). \texttt{r2} is ready in word 4, so
\texttt{add} cannot go earlier; \texttt{addi} shares word 5 with \texttt{sw} because
\texttt{sw} reads the old \texttt{r10}. A superscalar would find the same
parallelism at run time in hardware, with no \texttt{nop}s stored in the code, and the
same binary would still run on a wider or differently timed implementation.
\end{noot}

\begin{noot}{Exercise~\ref{q:mp-sync}}
\begin{enumerate}[label=(\alph*),leftmargin=*]
  \item While the lock is held, \texttt{ll}/\texttt{bne} just reads the cached copy
  (state S), so spinning causes no bus traffic. Only when the owner releases the lock
  (an invalidation) do the waiters miss, see 0 and try \texttt{sc}. A tight loop of
  atomic exchanges writes on every attempt, so the block ping-pongs between caches in
  M state with an invalidation on each try, flooding the bus and slowing the owner.
  \item The last thread resets the counter and sets \texttt{release=1}. A fast
  thread may leave and reach the barrier again, resetting \texttt{release=0} before a
  slow thread has seen the 1; the slow thread spins forever, and the fast threads
  wait for it: deadlock. With sense reversal each episode waits for the opposite
  value of \texttt{release}, so the next episode never undoes the previous release.
\end{enumerate}
\end{noot}

\begin{noot}{Exercise~\ref{q:mp-misc}}
\begin{enumerate}[label=(\alph*),leftmargin=*]
  \item $0.8\times20+0.2\times100=36$~ns; $0.5\times20+0.5\times100=60$~ns. Data
  placement matters on NUMA.
  \item Thread A's long misses: coarse-grained switching (or any style). Thread B's
  short stalls: only fine-grained or SMT, since coarse-grained switching costs a
  pipeline refill that is longer than the stall.
  \item The branch is resolved in a stage further from IF, so more stages (cycles)
  of wrongly fetched instructions must be flushed, even though each cycle is
  shorter.
\end{enumerate}
\end{noot}



\end{document}
